\documentclass[12pt,a4paper]{article}

% -------------------------------------------------------------------
% Codificação, idioma e fonte
% -------------------------------------------------------------------
\usepackage[T1]{fontenc}
\usepackage[utf8]{inputenc}
\usepackage[brazil]{babel}
\usepackage{lmodern}
\usepackage{tikz}
\usepackage{pgfplots}
\pgfplotsset{compat=1.18}
\usepackage{xparse}
\usepackage{float}
\usepackage[charter]{mathdesign}

% -------------------------------------------------------------------
% Cores
% -------------------------------------------------------------------
\usepackage[dvipsnames]{xcolor}
\definecolor{equilibriocolor}{RGB}{190,40,40}
\definecolor{taxacolor}{RGB}{0,90,160}
\definecolor{solucaoum}{RGB}{0,114,178}
\definecolor{solucaodois}{RGB}{126,47,142}
\definecolor{solucaotres}{RGB}{0,140,90}
\definecolor{solucaoquatro}{RGB}{230,159,0}
\definecolor{particularcolor}{RGB}{20,20,20}
\definecolor{alert}{RGB}{200,0,0}
\definecolor{structure}{rgb}{0.2,0.2,0.7}
\colorlet{example}{green!50!black}
\usepackage[ruled,vlined,linesnumbered]{algorithm2e}

\SetKwInput{Entrada}{Entrada}
\SetKwInput{Saida}{Saída}
\SetKwBlock{Inicio}{início}{fim}
\SetKwFor{ParaCada}{para cada}{faça}{fim}
\SetKwFor{Para}{para}{faça}{fim}
\SetKw{Retorna}{retorna}
% -------------------------------------------------------------------
% Layout da página
% -------------------------------------------------------------------
\usepackage[a4paper,margin=2.5cm]{geometry}
\setlength{\parindent}{0pt}
\setlength{\parskip}{0.6em}

% -------------------------------------------------------------------
% Matemática e símbolos
% -------------------------------------------------------------------

\usepackage{amsmath,amssymb}
\usepackage{cancel}
\usepackage{mathtools}
\usepackage{array}
\usepackage{microtype}
\usepackage{enumitem}
\usepackage{nicematrix}
\usepackage{graphicx}
% -------------------------------------------------------------------
% Gráficos, caixas e figuras
% -------------------------------------------------------------------
\usepackage[most]{tcolorbox}

% Versão compatível fora do beamer:
% \alt<overlay>{texto A}{texto B} -> mostra sempre texto A
\NewDocumentCommand{\alt}{d<> m m}{#2}

\newtcolorbox{block}[1][]{
  enhanced,
  colback=blue!5,
  colframe=structure,
  fonttitle=\bfseries,
  title={#1},
  rounded corners,
  boxrule=0.8pt,
  drop shadow
}

% -------------------------------------------------------------------
% Qualidade tipográfica
% -------------------------------------------------------------------
\usepackage{microtype}

% -------------------------------------------------------------------
% Listas
% -------------------------------------------------------------------
\usepackage{enumitem}
\setlist[itemize]{leftmargin=2em}
\setlist[enumerate]{leftmargin=2.2em}

% -------------------------------------------------------------------
% Legendas e links
% -------------------------------------------------------------------
\usepackage{caption}
\captionsetup{labelfont=bf,font=small}

\usepackage{hyperref}
\hypersetup{
  colorlinks=true,
  linkcolor=MidnightBlue,
  urlcolor=MidnightBlue,
  citecolor=MidnightBlue
}

% -------------------------------------------------------------------
% Cabeçalho e rodapé
% -------------------------------------------------------------------
\usepackage{fancyhdr}
\pagestyle{fancy}
\fancyhf{}
\lhead{Lista 1 - Modelo 2SLS}
\rhead{Economia Industrial}
\cfoot{\thepage}
\renewcommand{\headrulewidth}{0.4pt}
\setlength{\headheight}{14.5pt}

% -------------------------------------------------------------------
% Ambiente "conditions"
% -------------------------------------------------------------------
\newenvironment{conditions}
  {\par\vspace{\abovedisplayskip}\noindent\begin{tabular}{>{$}l<{$} @{${}={}$} l}}
  {\end{tabular}\par\vspace{\belowdisplayskip}}

% -------------------------------------------------------------------
% Título
% -------------------------------------------------------------------
% -------------------------------------------------------------------
% Título
% -------------------------------------------------------------------
\title{\textbf{Lista 1 - Economia Industrial - Modelo 2SLS}}

\author{
\begin{tabular}{ccc}
\begin{tabular}{c}
Luiz Mario Andrade\\
\small (Matrícula: 252029360)
\end{tabular}
&
\begin{tabular}{c}
Felipe Santos\\
\small (Matrícula: 232010719)
\end{tabular}
&
\begin{tabular}{c}
Luiza Nodari\\
\small (Matrícula: 242011335)
\end{tabular}
\\[1.2em]
\multicolumn{3}{c}{
\begin{tabular}{c}
Diogo Martins\\
\small (Matrícula: 232001578)
\end{tabular}
\begin{tabular}{c}
Sarah Moura\\
\small (Matrícula: 211060316)
\end{tabular}
\begin{tabular}{c}
Pedro Bijos\\
\small (Matrícula: 241003849)
\end{tabular}
}
\end{tabular}
}

\date{\today}
% -------------------------------------------------------------------
% Documento
% -------------------------------------------------------------------
\begin{document}

\maketitle

%\tableofcontents
\bigskip

% ================================================================
% ================================================================
\section{Introdução}

Esta atividade consiste em uma aplicação direta da estimação de demanda por um produto homogêneo, tomando como objeto empírico o mercado de frango. A base utilizada é a mesma do exemplo de demanda agregada por galeto discutido nas notas de aula, inspirado em Epple e McCallum e Reiss, no qual se estima uma demanda log-log para um produto homogêneo usando quantidade consumida, preço do próprio produto, renda e preço de um substituto. Aqui, adotamos a notação da lista de exercícios e interpretamos o produto como frango.

O ponto de partida é a especificação estrutural da função de demanda:

\[
q_t = \beta_0 + \beta_p p^{ch}_t + \beta_y y_t + \beta_b p^b_t + u_t,
\tag{1}
\label{eq:demanda_estrutural}
\]

em que \(q_t\) representa o logaritmo natural da quantidade consumida per capita de frango, \(p^{ch}_t\) representa o logaritmo natural do preço real do frango, \(y_t\) representa o logaritmo natural da renda real per capita e \(p^b_t\) representa o logaritmo natural do preço real da carne bovina, tratada como bem substituto. O parâmetro central da análise é \(\beta_p\), interpretado como a elasticidade-preço própria da demanda.

A formulação log-log permite interpretar diretamente os coeficientes como elasticidades. Assim,

\[
\beta_p
=
\frac{\partial q_t}{\partial p^{ch}_t}
=
\frac{\partial \ln Q_t}{\partial \ln P^{ch}_t},
\]

de modo que \(\beta_p\) mede a variação percentual esperada na quantidade demandada de frango diante de uma variação de 1\% no preço real do frango, mantendo constantes a renda e o preço da carne bovina. Portanto, se \(|\beta_p|<1\), a demanda é inelástica; se \(|\beta_p|>1\), a demanda é elástica.

Do ponto de vista de Economia Industrial, essa elasticidade é fundamental porque está ligada à análise de poder de mercado, margem preço-custo e perda crítica. Em uma formulação simples de firma com poder de mercado, a regra de elasticidade pode ser representada por:

\[
\frac{P - CMg}{P}
=
-\frac{1}{\varepsilon_p},
\tag{2}
\label{eq:lerner}
\]

em que \(P\) é o preço, \(CMg\) é o custo marginal e \(\varepsilon_p\) é a elasticidade-preço da demanda. Quanto menor a elasticidade em módulo, maior tende a ser a margem compatível com o equilíbrio. Em análises antitruste, a elasticidade também entra em exercícios de perda crítica. Para um aumento percentual de preço \(\Delta p\) e uma margem pré-aumento \(m\), uma expressão usual da perda crítica é:

\[
CL
=
\frac{\Delta p}{\Delta p + m},
\tag{3}
\label{eq:critical_loss}
\]

em que \(CL\) representa a fração de vendas que poderia ser perdida antes que o aumento de preço deixasse de ser lucrativo.

A principal dificuldade econométrica da estimação da equação~\eqref{eq:demanda_estrutural} é a potencial endogeneidade do preço do frango. Como preço e quantidade são determinados conjuntamente no equilíbrio de mercado, choques não observados de demanda, capturados por \(u_t\), podem afetar simultaneamente \(q_t\) e \(p^{ch}_t\). Nesse caso,

\[
\operatorname{Cov}(p^{ch}_t,u_t)\neq 0,
\tag{4}
\label{eq:endogeneidade}
\]

violando a condição necessária para que o MQO identifique causalmente \(\beta_p\). O viés do MQO pode ser representado, de forma simplificada, por:

\[
\operatorname{plim}\hat{\beta}^{MQO}_p
=
\beta_p
+
\frac{\operatorname{Cov}(p^{ch}_t,u_t)}
{\operatorname{Var}(p^{ch}_t)}.
\tag{5}
\label{eq:vies_mqo}
\]

Se choques positivos de demanda elevam simultaneamente quantidade e preço, então a covariância entre preço e erro tende a ser positiva. Como a verdadeira elasticidade-preço da demanda é negativa, esse viés positivo empurra a estimativa de MQO em direção a zero, fazendo a demanda parecer mais inelástica do que realmente é.

Para lidar com a endogeneidade, utiliza-se o método de variáveis instrumentais. O instrumento principal é o preço real do milho, denotado por \(z_t\), que atua como deslocador de custos no setor de frango. A hipótese econômica é que o milho afeta o preço do frango pelo lado da oferta, mas não afeta diretamente a demanda final por frango, condicionalmente à renda e ao preço da carne bovina.

As duas condições fundamentais para validade do instrumento são:

\[
\operatorname{Cov}(z_t,p^{ch}_t)\neq 0
\tag{6}
\label{eq:relevancia}
\]

e

\[
\operatorname{Cov}(z_t,u_t)=0.
\tag{7}
\label{eq:exogeneidade}
\]

A primeira é a condição de relevância: o instrumento precisa explicar a variável endógena. A segunda é a condição de exogeneidade ou exclusão: o instrumento não pode estar correlacionado com o erro estrutural da demanda.

O primeiro estágio do modelo de variáveis instrumentais é dado por:

\[
p^{ch}_t
=
\pi_0 + \pi_z z_t + \pi_y y_t + \pi_b p^b_t + v_t.
\tag{8}
\label{eq:primeiro_estagio}
\]

A hipótese de relevância é testada por:

\[
H_0:\pi_z=0.
\]

No caso de um único instrumento excluído, a força do instrumento pode ser avaliada pelo teste \(F\) do primeiro estágio. Em termos gerais, para \(q\) instrumentos excluídos, o teste \(F\) usual pode ser escrito como:

\[
F
=
\frac{(RSS_R - RSS_U)/q}{RSS_U/(n-k)},
\tag{9}
\label{eq:f_usual}
\]

em que \(RSS_R\) é a soma dos quadrados dos resíduos do modelo restrito, \(RSS_U\) é a soma dos quadrados dos resíduos do modelo irrestrito, \(q\) é o número de restrições testadas, \(n\) é o número de observações e \(k\) é o número de parâmetros no primeiro estágio. A regra de bolso tradicional considera que \(F>10\) reduz a preocupação com instrumentos fracos, embora essa regra possa ser inadequada sob heterocedasticidade.

O estimador de Mínimos Quadrados em Dois Estágios, ou 2SLS, pode ser representado em forma matricial por:

\[
\hat{\beta}_{2SLS}
=
(X'P_ZX)^{-1}X'P_ZY,
\tag{10}
\label{eq:2sls}
\]

em que \(P_Z = Z(Z'Z)^{-1}Z'\) é a matriz de projeção sobre o espaço gerado pelos instrumentos. Intuitivamente, o 2SLS substitui a variável endógena por sua parte prevista pelos instrumentos e pelos controles exógenos.

No caso de uma única variável endógena e um único instrumento excluído, a lógica do IV pode ser apresentada pela razão entre forma reduzida e primeiro estágio:

\[
\hat{\beta}^{IV}_p
=
\frac{\widehat{\operatorname{Cov}}(z_t,q_t)}
{\widehat{\operatorname{Cov}}(z_t,p^{ch}_t)}
=
\frac{\hat{\delta}}{\hat{\gamma}},
\tag{11}
\label{eq:iv_ratio}
\]

em que \(\hat{\delta}\) representa a forma reduzida da quantidade em relação ao instrumento e \(\hat{\gamma}\) representa o primeiro estágio do preço em relação ao instrumento. Essa expressão ajuda a entender o problema de instrumentos fracos: quando \(\hat{\gamma}\) se aproxima de zero, o denominador da razão fica pequeno e a estimativa de \(\beta_p\) se torna instável.

Além do 2SLS, também se utiliza o Método Generalizado dos Momentos, ou GMM. A condição de momento básica é:

\[
\mathbb{E}[Z_t u_t(\beta)] = 0,
\tag{12}
\label{eq:momento_gmm}
\]

e sua contraparte amostral é:

\[
g_n(\beta)
=
\frac{1}{n}\sum_{t=1}^{n} Z_t u_t(\beta).
\tag{13}
\label{eq:momento_amostral}
\]

O estimador GMM escolhe o vetor de parâmetros que minimiza a forma quadrática:

\[
\hat{\beta}_{GMM}
=
\arg\min_{\beta}
\left[
g_n(\beta)' W_n g_n(\beta)
\right],
\tag{14}
\label{eq:gmm_objetivo}
\]

em que \(W_n\) é uma matriz de ponderação positiva definida. No caso exatamente identificado, o número de momentos é igual ao número de parâmetros a identificar, de modo que a escolha da matriz de ponderação não altera a estimativa pontual. No caso sobreidentificado, há mais momentos do que parâmetros, e a matriz de ponderação afeta a combinação ótima das condições de momento.

A validade conjunta das restrições sobreidentificadoras é avaliada pelo teste \(J\) de Hansen:

\[
J
=
n \cdot g_n(\hat{\beta}_{GMM})'
\hat{W}
g_n(\hat{\beta}_{GMM}),
\tag{15}
\label{eq:hansen_j}
\]

com distribuição assintótica:

\[
J \sim \chi^2_{m-k},
\tag{16}
\label{eq:hansen_dist}
\]

em que \(m\) é o número de instrumentos e \(k\) é o número de variáveis endógenas. A hipótese nula do teste é a validade conjunta das restrições sobreidentificadoras. Rejeitar a hipótese nula sugere que ao menos uma condição de momento pode estar violada.

Também são discutidos instrumentos fracos. Um instrumento é fraco quando possui baixa correlação com a variável endógena, mesmo que seja exógeno. Nesse caso, o 2SLS pode ser viesado, altamente disperso e produzir intervalos de confiança convencionais pouco confiáveis. Por isso, além do \(F\) usual, utiliza-se o \(F\) efetivo de Montiel Olea e Pflueger, denotado por:

\[
F^{Ef},
\tag{17}
\label{eq:f_efetivo}
\]

que compara a força do primeiro estágio com valores críticos robustos à heterocedasticidade. A regra de decisão geral é:

\[
F^{Ef} > CV_{\tau}
\quad \Rightarrow \quad
\text{rejeita-se a hipótese de instrumento fraco para o nível de tolerância } \tau.
\tag{18}
\label{eq:regra_mop}
\]

Quando a identificação é fraca, também podem ser usados intervalos robustos, como Anderson--Rubin ou CLR. A lógica do teste Anderson--Rubin é testar valores candidatos \(\beta_0\) por meio da regressão:

\[
q_t - \beta_0 p^{ch}_t
=
\alpha_0 + \alpha_y y_t + \alpha_b p^b_t + Z_t'\alpha_z + e_t,
\tag{19}
\label{eq:ar}
\]

avaliando se os instrumentos são conjuntamente significativos nessa equação. O conjunto de valores de \(\beta_0\) não rejeitados forma o intervalo de confiança robusto a instrumentos fracos.

Por fim, a lista utiliza uma simulação para ilustrar a deterioração do estimador 2SLS quando a força do instrumento diminui. O desenho básico da simulação pode ser representado por:

\[
x_i = \pi z_i + v_i,
\tag{20}
\label{eq:sim_primeiro_estagio}
\]

\[
y_i = \beta x_i + u_i,
\tag{21}
\label{eq:sim_estrutural}
\]

com \(\operatorname{Cov}(u_i,v_i)\neq 0\) para gerar endogeneidade. À medida que \(\pi\) diminui, o instrumento perde força, o primeiro estágio se enfraquece e o 2SLS passa a apresentar maior viés, maior dispersão e pior comportamento inferencial.

Assim, a atividade articula conceitos fundamentais de Economia Industrial e Econometria Aplicada: estimação de demanda, identificação da elasticidade-preço, endogeneidade de preços, variáveis instrumentais, 2SLS, GMM, sobreidentificação, teste \(J\) de Hansen, instrumentos fracos, inferência robusta e implicações antitruste da elasticidade estimada.

% ================================================================
\section{Base de Dados e Tratamento}

A base de dados utilizada é a mesma do exemplo de demanda agregada por galeto apresentado nas notas de aula, adaptada à notação da lista para demanda por frango. Trata-se de uma base anual com informações de quantidade consumida, renda, preços de frango, carne bovina, milho, ração, índice de preços ao consumidor, produção agregada, população, exportações de carnes e tendência temporal.

As variáveis originais disponíveis na base são:

\begin{conditions}
Q        & \text{consumo per capita de frango, em libras, equivalente desossado;}\\
Y        & \text{renda disponível real per capita;}\\
PCHICK   & \text{índice de preços ao consumidor para frango;}\\
PBEEF    & \text{índice de preços ao consumidor para carne bovina;}\\
PCOR     & \text{índice de preços ao produtor para milho;}\\
PF       & \text{índice nominal de preço da ração de frango;}\\
CPI      & \text{índice de preços ao consumidor;}\\
QPROD    & \text{produção agregada de frangos jovens;}\\
POP      & \text{população dos Estados Unidos em 1º de julho;}\\
EATEX    & \text{exportações de carne bovina, vitela e suína;}\\
TIM      & \text{tendência temporal usada nas regressões.}
\end{conditions}

A preparação dos dados consistiu em deflacionar os preços nominais pelo índice de preços ao consumidor e aplicar logaritmos naturais às variáveis utilizadas na demanda. As principais variáveis construídas foram:

\[
q_t = \ln(Q_t),
\tag{22}
\label{eq:q_log}
\]

\[
y_t = \ln(Y_t),
\tag{23}
\label{eq:y_log}
\]

\[
p^{ch}_t
=
\ln\left(\frac{PCHICK_t}{CPI_t}\right),
\tag{24}
\label{eq:pch_log}
\]

\[
p^b_t
=
\ln\left(\frac{PBEEF_t}{CPI_t}\right),
\tag{25}
\label{eq:pb_log}
\]

\[
z_t
=
\ln\left(\frac{PCOR_t}{CPI_t}\right).
\tag{26}
\label{eq:z_log}
\]

Assim, \(p^{ch}_t\) mede o logaritmo do preço real do frango, \(p^b_t\) mede o logaritmo do preço real da carne bovina e \(z_t\) mede o logaritmo do preço real do milho. O termo \(z_t\) é utilizado como instrumento principal para o preço real do frango, pois o milho é um insumo relevante na produção de frango e pode deslocar custos de produção.

Para as especificações alternativas de instrumentos, foram construídas potências e defasagens do instrumento:

\[
z_t^2 =
\left[
\ln\left(\frac{PCOR_t}{CPI_t}\right)
\right]^2,
\tag{27}
\label{eq:z2}
\]

\[
z_t^3 =
\left[
\ln\left(\frac{PCOR_t}{CPI_t}\right)
\right]^3,
\tag{28}
\label{eq:z3}
\]

\[
z_{t-1}
=
L.z_t,
\qquad
z_{t-1}^2
=
(L.z_t)^2.
\tag{29}
\label{eq:z_lags}
\]

Os conjuntos de instrumentos analisados foram:

\[
\begin{aligned}
Z_1 &= \{z_t\},\\
Z_2 &= \{z_t,z_t^2\},\\
Z_3 &= \{z_t,z_t^2,z_t^3\},\\
Z_4 &= \{z_{t-1}\},\\
Z_5 &= \{z_{t-1},z_{t-1}^2\},\\
Z_6 &= \{z_t,z_{t-1}\},\\
Z_7 &= \{z_t,z_t^2,z_{t-1},z_{t-1}^2\}.
\end{aligned}
\tag{30}
\label{eq:conjuntos_instrumentos}
\]

A equação estrutural estimada em todas as especificações mantém a mesma forma:

\[
q_t = \beta_0 + \beta_p p^{ch}_t + \beta_y y_t + \beta_b p^b_t + u_t.
\tag{31}
\label{eq:base_estrutural}
\]

A diferença entre as especificações está no conjunto de instrumentos excluídos utilizado para identificar o preço real do frango. O primeiro estágio geral pode ser escrito como:

\[
p^{ch}_t
=
\pi_0
+
Z_t'\pi_z
+
\pi_y y_t
+
\pi_b p^b_t
+
v_t,
\tag{32}
\label{eq:base_primeiro_estagio}
\]

em que \(Z_t\) representa um dos conjuntos de instrumentos definidos em \eqref{eq:conjuntos_instrumentos}.

O diagnóstico de relevância dos instrumentos foi baseado no \(R^2\) parcial e no teste \(F\) do primeiro estágio. O \(R^2\) parcial do instrumento excluído pode ser interpretado como a parcela da variação de \(p^{ch}_t\) explicada pelos instrumentos, após controlar por \(y_t\) e \(p^b_t\). Uma forma usual de expressá-lo é:

\[
R^2_{\text{parcial}}
=
\frac{R^2_{\text{irrestrito}} - R^2_{\text{restrito}}}
{1 - R^2_{\text{restrito}}},
\tag{33}
\label{eq:r2_parcial}
\]

em que o modelo irrestrito inclui os instrumentos excluídos e os controles, enquanto o modelo restrito inclui apenas os controles.

Para os modelos sobreidentificados, a validade conjunta dos instrumentos foi avaliada pelo teste \(J\) de Hansen:

\[
H_0:
\mathbb{E}[Z_tu_t]=0.
\tag{34}
\label{eq:h0_hansen}
\]

A rejeição dessa hipótese indica que o conjunto de instrumentos pode não ser conjuntamente exógeno. Já a não rejeição não prova que todos os instrumentos sejam válidos individualmente, mas apenas que não há evidência estatística suficiente contra as restrições sobreidentificadoras.

A Figura~\ref{fig:series_log} apresenta as séries logarítmicas utilizadas no exercício. Como a especificação é log-log, a escala logarítmica é coerente com a interpretação dos coeficientes como elasticidades. A figura também permite visualizar a evolução conjunta da quantidade consumida, dos preços reais e da renda ao longo do tempo, ajudando a contextualizar a estimação da demanda.

\begin{figure}[H]
    \centering
    \includegraphics[width=0.90\textwidth]{Lista_1/Figures/Stata/stata_fig01_series.pdf}
    \caption{Séries logarítmicas utilizadas na estimação da demanda por frango. A figura apresenta a evolução temporal da quantidade consumida per capita, do preço real do frango, da renda real per capita, do preço real da carne bovina e do preço real do milho.}
    \label{fig:series_log}
\end{figure}
\section{Implementação Computacional}

A implementação computacional foi realizada nos softwares \textbf{Stata} e \textbf{R}. O Stata foi utilizado como ambiente principal de estimação e replicação, enquanto os códigos em R foram estruturados como uma tradução da rotina em Stata, reproduzindo a mesma sequência lógica de tratamento da base, estimação dos modelos, construção das tabelas e geração das figuras. Em particular, o Stata foi mantido como referência para o teste $F$ efetivo de Montiel Olea e Pflueger, implementado por meio do comando \texttt{weakivtest}.

A execução principal no Stata é feita pelo arquivo:

\begin{verbatim}
Stata/00_master.do
\end{verbatim}

Esse script funciona como arquivo mestre da replicação. Ele carrega as configurações gerais, define os diretórios do projeto, abre o log de execução e chama sequencialmente os demais scripts. A estrutura geral do projeto foi organizada da seguinte forma:

\begin{verbatim}
lista1EI_26_replication_package/
├── data/
│   ├── raw/                 # bases originais: chicken.dta e chicken.csv
│   └── processed/           # bases tratadas
├── output/
│   ├── tables/              # tabelas em .csv, .dta e .tex
│   ├── figures/             # figuras em .png e .pdf
│   └── logs/                # logs de execução
├── Stata/                   # scripts .do
└── R/                       # scripts .R
\end{verbatim}

Os caminhos são definidos em \texttt{Stata/config.do}, que cria os globais \texttt{ROOT}, \texttt{RAW}, \texttt{PROC}, \texttt{TABS}, \texttt{FIGS} e \texttt{LOGS}. Assim, as bases originais são lidas a partir de \texttt{data/raw}, as bases tratadas são salvas em \texttt{data/processed}, as tabelas em \texttt{output/tables}, as figuras em \texttt{output/figures} e os logs em \texttt{output/logs}. Para permitir que cada script seja executado isoladamente, também foi criado o arquivo auxiliar:

\begin{verbatim}
Stata/_load_config_if_needed.do
\end{verbatim}

Esse arquivo verifica se os caminhos globais já foram carregados e, caso contrário, executa automaticamente o \texttt{config.do}.

\subsection*{Tratamento da base}

O tratamento da base é realizado pelo script:

\begin{verbatim}
Stata/01_prepare_data.do
\end{verbatim}

Esse script importa a base original do exemplo de demanda por galeto/frango, preferencialmente em formato \texttt{.dta}. Caso o arquivo \texttt{chicken.dta} não esteja disponível, o código importa a versão em \texttt{.csv}. Em seguida, são construídas as variáveis em logaritmo definidas nas equações~\eqref{eq:q_log}, \eqref{eq:y_log}, \eqref{eq:pch_log}, \eqref{eq:pb_log} e \eqref{eq:z_log}. Também são criadas as potências e defasagens do instrumento, conforme as equações~\eqref{eq:z2}, \eqref{eq:z3} e \eqref{eq:z_lags}.

A base tratada é salva em:

\begin{verbatim}
data/processed/chicken_prepared_stata.dta
data/processed/chicken_prepared_stata.csv
\end{verbatim}

Esses arquivos são usados como entrada pelos demais scripts de estimação.

\subsection*{Estimações de MQO, 2SLS e primeiro estágio}

As estimações iniciais são realizadas pelo script:

\begin{verbatim}
Stata/02_ols_iv_first_stage.do
\end{verbatim}

Primeiro, o código estima por MQO a equação estrutural de demanda definida em~\eqref{eq:demanda_estrutural}, utilizando erros-padrão robustos à heterocedasticidade. Os resultados são salvos em:

\begin{verbatim}
output/tables/stata_question_01_ols.dta
output/tables/stata_question_01_ols.csv
\end{verbatim}

Em seguida, o script estima o modelo por 2SLS usando \(Z_1=\{z_t\}\) como instrumento excluído para \(p_t^{ch}\). O procedimento computacional segue a lógica do primeiro estágio apresentado em~\eqref{eq:primeiro_estagio} e do estimador 2SLS apresentado em~\eqref{eq:2sls}. Os resultados da estimação 2SLS são salvos em:

\begin{verbatim}
output/tables/stata_question_03_iv_z1.dta
output/tables/stata_question_03_iv_z1.csv
\end{verbatim}

O mesmo script calcula o primeiro estágio, o \(R^2\) parcial do instrumento excluído e a estatística \(F\) usual associada ao teste de relevância. O cálculo do \(R^2\) parcial segue a expressão~\eqref{eq:r2_parcial}. A tabela do primeiro estágio é salva em:

\begin{verbatim}
output/tables/stata_question_04_first_stage.dta
output/tables/stata_question_04_first_stage.csv
\end{verbatim}

Além disso, o script gera uma base auxiliar com resíduos e valores ajustados do MQO, usada posteriormente nos gráficos de diagnóstico:

\begin{verbatim}
data/processed/chicken_with_ols_residuals_stata.dta
\end{verbatim}

\subsection*{GMM e teste de Hansen}

A estimação por GMM é realizada pelo script:

\begin{verbatim}
Stata/03_gmm_hansen.do
\end{verbatim}

Esse script estima os modelos GMM com os conjuntos \(Z_1=\{z_t\}\) e \(Z_2=\{z_t,z_t^2\}\). A implementação segue as condições de momento definidas em~\eqref{eq:momento_gmm} e~\eqref{eq:momento_amostral}, bem como a função objetivo do GMM apresentada em~\eqref{eq:gmm_objetivo}.

No modelo exatamente identificado, \(Z_1\), a escolha da matriz de ponderação não altera a estimativa pontual, pois o número de momentos coincide com o número de parâmetros a identificar. No modelo sobreidentificado, \(Z_2\), o código calcula a estatística \(J\) de Hansen conforme~\eqref{eq:hansen_j}, usando a distribuição assintótica indicada em~\eqref{eq:hansen_dist}. Os resultados são salvos em:

\begin{verbatim}
output/tables/stata_question_05_gmm.dta
output/tables/stata_question_05_gmm.csv
\end{verbatim}

\subsection*{Especificações alternativas de instrumentos}

As especificações alternativas são estimadas pelo script:

\begin{verbatim}
Stata/04_alt_instruments_weakiv.do
\end{verbatim}

Esse arquivo estima todos os conjuntos de instrumentos definidos em~\eqref{eq:conjuntos_instrumentos}. Para cada especificação \(Z_j\), o algoritmo executa as seguintes etapas:

\begin{enumerate}[label=(\roman*)]
    \item estima o primeiro estágio geral, conforme~\eqref{eq:base_primeiro_estagio};
    \item calcula o \(R^2\) parcial, conforme~\eqref{eq:r2_parcial};
    \item calcula o teste \(F\) usual do primeiro estágio, conforme a lógica de~\eqref{eq:f_usual};
    \item estima a equação estrutural por 2SLS, conforme~\eqref{eq:2sls};
    \item calcula o teste $J$ de Hansen quando o modelo é sobreidentificado, conforme~\eqref{eq:hansen_j} e~\eqref{eq:hansen_dist};
    \item calcula, quando disponível, o $F$ efetivo de Montiel Olea e Pflueger, conforme a regra de decisão em~\eqref{eq:regra_mop}.
\end{enumerate}

As tabelas principais geradas por esse script são:

\begin{verbatim}
output/tables/stata_question_08_first_stage_all.dta
output/tables/stata_question_08_first_stage_all.csv
output/tables/stata_question_09_comparative.dta
output/tables/stata_question_09_comparative.csv
\end{verbatim}

Além disso, para a questão sobre validade dos instrumentos, o código gera a tabela específica do teste \(J\) de Hansen para modelos sobreidentificados:

\begin{verbatim}
output/tables/stata_question_10_hansen_overid.csv
output/tables/stata_question_10_hansen_overid.tex
\end{verbatim}

e o gráfico correspondente dos valores-$p$:

\begin{verbatim}
output/figures/stata_fig12_hansen_pvalues.png
output/figures/stata_fig12_hansen_pvalues.pdf
\end{verbatim}

\subsection*{Inferência robusta a instrumentos fracos}

Os intervalos robustos a instrumentos fracos são calculados pelo script:

\begin{verbatim}
Stata/05_ar_intervals.do
\end{verbatim}

A rotina implementa uma busca em grade para valores candidatos de \(\beta_p\). Para cada valor candidato, o código aplica a transformação associada ao teste de Anderson--Rubin descrita em~\eqref{eq:ar}. Em seguida, verifica se os instrumentos excluídos são conjuntamente significativos na regressão auxiliar. Os valores de \(\beta_p\) para os quais a hipótese não é rejeitada compõem o intervalo robusto.

Como instrumentos fracos podem gerar intervalos abertos, muito amplos ou disjuntos, o código utiliza uma grade adaptativa para capturar esses formatos. A tabela final é salva em:

\begin{verbatim}
output/tables/stata_question_11_ar_intervals.dta
output/tables/stata_question_11_ar_intervals.csv
\end{verbatim}

\subsection*{Gráficos}

Os gráficos principais são gerados pelo script:

\begin{verbatim}
Stata/06_visualizations.do
\end{verbatim}

Esse arquivo utiliza a base tratada e as tabelas produzidas nas etapas anteriores para construir as figuras usadas no texto. As principais saídas são:

\begin{verbatim}
output/figures/stata_fig01_series.pdf
output/figures/stata_fig02_scatter_demand.pdf
output/figures/stata_fig03_residuals_fitted.pdf
output/figures/stata_fig04_beta_p_mqo_versus_2sls.pdf
output/figures/stata_fig11_gmm_beta_hansen.pdf
\end{verbatim}

A Figura~\ref{fig:series_log} é construída a partir das variáveis definidas nas equações~\eqref{eq:q_log}, \eqref{eq:y_log}, \eqref{eq:pch_log}, \eqref{eq:pb_log} e~\eqref{eq:z_log}. A figura de comparação entre MQO e 2SLS usa os resultados da estimação da equação~\eqref{eq:demanda_estrutural} por MQO e por 2SLS, conforme~\eqref{eq:2sls}. Já o gráfico de diagnóstico dos resíduos utiliza a base auxiliar com resíduos do MQO.

O gráfico de instabilidade IV com os dados da lista é gerado por:

\begin{verbatim}
Stata/08_iv_curve_actual_data.do
\end{verbatim}

Esse script operacionaliza a relação apresentada em~\eqref{eq:iv_ratio}. O código estima a forma reduzida e o primeiro estágio nos dados, mantém \(\hat{\delta}\) fixo e varia contrafactualmente \(\gamma\), mostrando como \(\hat{\beta}^{IV}_p=\hat{\delta}/\hat{\gamma}\) se torna instável quando o primeiro estágio se aproxima de zero. As saídas são:

\begin{verbatim}
output/figures/stata_fig10_iv_curve_actual_data.png
output/figures/stata_fig10_iv_curve_actual_data.pdf
\end{verbatim}

\subsection*{Simulação de instrumentos fracos}

A simulação da questão 14 é implementada pelo script:

\begin{verbatim}
Stata/07_simulation_weak_instruments.do
\end{verbatim}

O algoritmo segue o desenho apresentado nas equações~\eqref{eq:sim_primeiro_estagio} e~\eqref{eq:sim_estrutural}. A ideia é controlar diretamente a força do instrumento por meio do parâmetro \(\pi\). Foram usados os valores:

\[
\pi \in \{1,\ 0{,}5,\ 0{,}25,\ 0{,}10,\ 0{,}05,\ 0{,}02\}.
\]

Para cada valor de \(\pi\), o código executa \(250\) repetições com \(n=200\) observações. Em cada repetição, o algoritmo:

\begin{enumerate}[label=(\roman*)]
    \item gera um instrumento exógeno \(z_i\);
    \item gera dois choques aleatórios independentes;
    \item constrói \(u_i\) e \(v_i\) de forma correlacionada, induzindo endogeneidade entre a variável endógena e o erro estrutural;
    \item constrói a variável endógena segundo~\eqref{eq:sim_primeiro_estagio};
    \item constrói a variável dependente segundo~\eqref{eq:sim_estrutural};
    \item estima o modelo por MQO;
    \item estima o primeiro estágio e armazena o \(F\) usual;
    \item estima o modelo por 2SLS;
    \item calcula o intervalo convencional de 95\%;
    \item registra se o intervalo contém o verdadeiro valor \(\beta_p=-1\).
\end{enumerate}

Ao final das repetições, o código calcula, para cada valor de \(\pi\), a média das estimativas 2SLS, o viés em relação ao valor verdadeiro, o desvio-padrão das estimativas, o \(F\) médio do primeiro estágio e a taxa de cobertura dos intervalos convencionais. A tabela da simulação é salva em:

\begin{verbatim}
output/tables/stata_question_14_simulation.dta
output/tables/stata_question_14_simulation.csv
\end{verbatim}

Os gráficos da simulação são salvos em:

\begin{verbatim}
output/figures/stata_fig06_simulation_sd.pdf
output/figures/stata_fig07_simulation_bias_convergence.png
output/figures/stata_fig08_simulation_ffs.pdf
output/figures/stata_fig09_q14_simulacao_instrumentos_fracos.pdf
\end{verbatim}

A lógica computacional da simulação é mostrar que, quando \(\pi\) diminui, o instrumento perde força, o primeiro estágio enfraquece e o estimador 2SLS passa a apresentar maior viés, maior dispersão e pior comportamento inferencial. Esse mecanismo é justamente o que a relação em~\eqref{eq:iv_ratio} antecipa: quando o primeiro estágio se aproxima de zero, a razão que define o estimador IV torna-se instável.

\subsection*{Implementação em R}

A implementação em R reproduz a estrutura dos códigos em Stata. Os scripts em R foram concebidos como uma tradução da rotina principal, mantendo a mesma ordem de operações: leitura e tratamento da base, estimação por MQO, 2SLS, GMM, cálculo do teste de Hansen, especificações alternativas de instrumentos, intervalos robustos, gráficos e simulação de instrumentos fracos.

As saídas em R seguem o mesmo padrão de organização, com tabelas salvas em \texttt{output/tables} e figuras em \texttt{output/figures}. Quando os resultados são comparáveis diretamente, os arquivos em R são salvos com prefixo próprio, permitindo comparação com as saídas do Stata sem sobrescrever os arquivos principais.

Portanto, a implementação computacional foi desenhada para garantir replicabilidade e rastreabilidade. Cada resultado apresentado no texto está associado a um script específico, a uma tabela em \texttt{output/tables} ou a uma figura em \texttt{output/figures}. O Stata fornece a referência principal das estimações e testes, enquanto o R funciona como tradução e verificação paralela da rotina empírica.
\section{Resultados e Discussão}

    
\begin{enumerate}[label=\textbf{\arabic*.}]

\item \textbf{MQO.} Estime a função de demanda~\eqref{eq:demanda_estrutural} por mínimos quadrados ordinários. Reporte $\hat{\beta}_p$, seu erro-padrão robusto à heterocedasticidade, intervalo de confiança de 95\% e interprete economicamente o sinal e a magnitude da elasticidade estimada.

\textbf{Resultados}

A estimação reportada nesta questão corresponde à função de demanda log-log definida na Equação~\eqref{eq:demanda_estrutural}, construída a partir das transformações da base descritas nas Equações~\eqref{eq:q_log}, \eqref{eq:pch_log}, \eqref{eq:y_log} e~\eqref{eq:pb_log}.

 Estimando a função de demanda por MQO com erros padrão robustos, obtivemos os seguintes resultados para a variável de preço do frango:

\begin{table}[H]
    \centering
    \caption{Resultados do MQO Robusto}
    \label{tab:mqo_robusto}
    \begin{tabular}{lcccccc}
        \hline
        Modelo & $\hat{\beta}_p$ & Erro-padrão & IC 95\% inf. & IC 95\% sup. & $N$  \\
        \hline
        MQO Robusto & -0.265 & 0.067 & -0.400 & -0.130 & 52 \\
        \hline
    \end{tabular}
\end{table}

    Dada uma função do tipo log-log, o parâmetro $\hat{\beta}_p$ representa diretamente a elasticidade-preço próprio da demanda por frango. O sinal negativo indica uma relação negativa do preço com a quantidade demandada. Como este valor é menor do 1 (em módulo), podemos dizer que a demanda por frango é \textbf{inelástica} em relação ao seu preço. 
    Em termos percentuais, um aumento de 1\% no preço do frango está associado a uma queda de apenas 0,265\% na quantidade per capita demandada de frango, \textit{ceteris paribus}. Isso pode ser verificado empiricamente pelo gráfico ~\ref{fig:demanda_observada}
\begin{figure}[H]
    \centering
    \includegraphics[width=0.85\textwidth]{Lista_1/Figures/Stata/stata_fig02_scatter_demand.pdf}
    \caption{Relação observada entre o log da quantidade per capita de frango e o log do preço real do frango. A inclinação negativa do ajuste linear é compatível com a intuição de uma curva de demanda decrescente, embora essa relação bruta não resolva o problema de endogeneidade do preço.}
    \label{fig:demanda_observada}
\end{figure}

\item \textbf{Endogeneidade do preço.} Explique por que $p_t^{ch}$ pode ser endógeno
    em uma equação de demanda. Discuta a direção provável do viés de MQO quando choques positivos de demanda elevam simultaneamente quantidade e preço.

\textbf{Resultados}

A discussão abaixo se relaciona diretamente à condição de endogeneidade indicada na Equação~\eqref{eq:endogeneidade} e à expressão do viés do MQO apresentada na Equação~\eqref{eq:vies_mqo}.

    Em um mercado competitivo, o preço ($p_t^{ch}$) e a quantidade não são determinados isoladamente, mas sim simultaneamente pelo cruzamento das forças de oferta e demanda. Em uma equação de demanda, o termo de erro estocástico ($u_t$) engloba todos os fatores não observados que afetam o consumo. Quando ocorre um choque positivo de demanda, a curva se desloca para a direita. Dado que a curva de oferta possui inclinação positiva, esse deslocamento eleva não apenas a quantidade transacionada, mas também empurra o preço de equilíbrio do mercado para cima. Essa dinâmica de determinação simultânea cria uma correlação positiva intrínseca entre o preço e o termo de erro ($Cov(p_t^{ch}, u_t) > 0$), o que viola a hipótese clássica de exogeneidade estrita e torna a variável de preço endógena.

    Ao estimar essa relação por Mínimos Quadrados Ordinários (MQO), o método falha em isolar o efeito puro do preço, capturando também a sujeira dessa simultaneidade. Matematicamente, a esperança do estimador de MQO é igual ao verdadeiro parâmetro ($\beta$) somado a um termo de viés dado pela razão $\frac{Cov(p_t^{ch}, u_t)}{Var(p_t^{ch})}$. Pela teoria econômica, sabemos que a verdadeira elasticidade-preço é negativa. Contudo, como o choque positivo de demanda eleva o preço simultaneamente, a correlação matemática entre eles é positiva, tornando o termo de viés globalmente positivo. Ao somarmos um viés positivo a um parâmetro verdadeiro negativo, o estimador de MQO é empurrado em direção a zero. Portanto, a direção provável do viés de MQO é positiva (viés de atenuação), o que subestima a verdadeira sensibilidade da demanda e faz a curva parecer mais inelástica do que realmente é.

    Essa dinâmica teórica é perfeitamente refletida nos resultados empíricos da estimação. Quando o modelo foi estimado via MQO, a elasticidade-preço do frango encontrada foi de apenas $-0.186$. No entanto, ao corrigir a endogeneidade através do método de Variáveis Instrumentais (2SLS) — utilizando um instrumento forte para isolar os choques de oferta —, o coeficiente revelou a verdadeira sensibilidade da demanda, passando para $-0.882$. A diferença substancial entre os dois valores evidencia, na prática, o viés positivo do MQO, que mascarava a real elasticidade do mercado ao confundir o deslocamento da curva de demanda com um movimento ao longo dela.

\item \textbf{2SLS com um instrumento.} Estime o modelo (1) por 2SLS, usando
    $z_t$ como instrumento excluído para $p_t^{ch}$. As variáveis $y_t$ e $p_t^b$
    devem ser incluídas tanto na equação estrutural quanto no primeiro estágio.
    Reporte $\hat{\beta}^{2SLS}_p$, erro-padrão robusto, intervalo de confiança e compare com a estimativa por MQO.\\

\textbf{Resultados}

A estimação por 2SLS utiliza a equação estrutural~\eqref{eq:demanda_estrutural}, o primeiro estágio~\eqref{eq:primeiro_estagio} e a forma matricial do estimador de variáveis instrumentais apresentada na Equação~\eqref{eq:2sls}.
\begin{table}[H]
    \centering
    \caption{Resultados do 2SLS com $Z_1 = \{z_t\}$}
    \label{tab:2sls_z1}
    \begin{tabular}{lcccccc}
        \hline
        Modelo & $\hat{\beta}_p$ & Erro-padrão & IC 95\% inf. & IC 95\% sup. & $N$ & $F$ 1º estágio \\
        \hline
        $Z_1$ & -1.659 & 0.991 & -3.600 & 0.283 & 40 & 3.245 \\
        \hline
    \end{tabular}
\end{table}

Enquanto o $\hat{\beta}_p$ estimado por MQO, no item 1, tem um coeficiente de $-0.186$, na estimação por 2SLS, em que o log do preço real do milho foi usado como instrumento, seu valor estimado foi de $-1.659$, de modo que uma demanda, inicialmente inelástica, torna-se elástica (|$-1.659$|>$1$), ou seja, os consumidores se tornam mais sensíveis às variações no preço do frango, quando estimado por 2SLS. Além disso, tal coeficiente implica que um aumento de 1\% no preço do frango representa uma diminuição de 1,659\% em sua quantidade demandada.

Contudo, com o intervalo de confiança estimado [$-3.600, 0.283$], a um nível de significância de 5\%, não é possível rejeitar a hipótese nula de insignificância estatística da variável no modelo, o que é reforçado pelo desvio-padrão elevado. Logo, variações no preço real do frango não têm efeito estatístico significativo sobre a quantidade demandada. A partir da estatística F (de primeiro estágio), relativamente baixa*, pode-se inferir que o o preço real do milho é um instrumento fraco.

No gráfico abaixo, é possível analisar visualmente as estimativas propostas para o modelo estimado por MQO e por 2SLS (Z1).

*Considerando a regra de decisão F>$10$ para instrumentos fortes.

\begin{figure}[H]
    \centering
    \includegraphics[width=0.90\textwidth]{Lista_1/Figures/Stata/stata_fig04_beta_p_mqo_versus_2sls.pdf}
    \caption{Elasticidade-preço estimada por MQO e por 2SLS para diferentes conjuntos de instrumentos. Os pontos representam as estimativas de $\beta_p$ e as barras horizontais representam os intervalos de confiança de 95\%. A linha vertical indica elasticidade igual a zero.}
    \label{fig:mqo_2sls}
\end{figure}

\item \textbf{Primeiro estágio.} Estime o primeiro estágio
    \[
        p_t^{ch} = \pi_0 + \pi_z z_t + \pi_y y_t + \pi_b p_t^b + v_t.
    \]
    Reporte o coeficiente de $z_t$, o $R^2$ parcial do instrumento excluído e o
    teste $F$ para a hipótese $H_0: \pi_z = 0$. Explique o que significa relevância
    do instrumento e por que relevância estatística não é suficiente para validade
    do instrumento.

\textbf{Resultados}

O modelo estimado nesta questão operacionaliza o primeiro estágio descrito na Equação~\eqref{eq:primeiro_estagio}. A relevância do instrumento se relaciona à condição da Equação~\eqref{eq:relevancia}, enquanto o cálculo do $R^2$ parcial remete à Equação~\eqref{eq:r2_parcial} e o teste $F$ usual à Equação~\eqref{eq:f_usual}.

Ao estimar o primeiro estágio, obteve-se a seguinte equação do preço do frango: 
 \begin{equation}
    p_t^{ch} = 7.7844 + 0.1088 z_t - 0.7850 y_t - 0.0160 p_t^{b}
    \label{eq:primeiro_estagio_estimado}
\end{equation}
    Pelos resultados do modelo descrito, o coeficiente de $z_t$ é de cerca de 0,1088, em que a proporção da variação do preço do frango explicada unicamente pelo preço do milho equivale a aproximadamente 4,13\%, isto é, o R² parcial do instrumento excluído. Isto posto, a relevância do instrumento diz respeito a sua capacidade de ser estatisticamente correlacionado com a variável endógena condicionado pelos demais controles, $Corr(Z_i,X_i) \neq 0$, representados pelas variáveis exógenas renda e preço da carne bovina. No entanto, o diagnóstico reportado pelo teste $F$ para a hipótese H\textsubscript{0}  feita refuta a validade do instrumento. Isso porque a estatística F de 2,883, sob a luz da literatura econométrica, é substancialmente abaixo de 10, valor de corte sugerido por Staiger e Stock (1997), apontando para um instrumento fraco. 

    Portanto, o instrumento fraco reproduz estimativas enviesadas por 2SLS e induz a conclusões não confiáveis. Contudo, o parâmetro da nossa análise não seria suficiente caso o instrumento satisfizesse à condição de relevância. O procedimento teórico exige que a variável ateste à condição de exogeneidade, dada por $Corr(X_i,\varepsilon_i) = 0$. Em outras palavras, devemos nos certificar quais canais afetam a quantidade demanda. Pelo nosso exemplo prático, o preço do milho deve estar estritamente correlacionado com o preço do frango. Uma vez impactado por choques não observados, concentrados no termo de erro, a hipótese de exogeneidade é violada. Em particular, para um modelo exatamente identificado como esse, não há elementos metodológicos que possibilitem determinar objetivamente o pressuposto, dependendo de uma fundamentação contextual para sua aprovação.

\item \textbf{GMM e sobreidentificação.} Estime o modelo por GMM usando
    dois conjuntos de instrumentos excluídos:
    \begin{enumerate}[label=(\alph*)]
        \item $\{z_t\}$;
        \item $\{z_t, z_t^2\}$.
    \end{enumerate}
    Explique por que, no caso exatamente identificado, a escolha da matriz de
    ponderação do GMM não deve alterar a estimativa pontual. No caso
    sobreidentificado, reporte a estatística $J$ de Hansen e interprete o teste de
    restrições sobreidentificadoras.

\textbf{Resultados}

A estimação por GMM parte das condições de momento das Equações~\eqref{eq:momento_gmm} e~\eqref{eq:momento_amostral}, e minimiza a função objetivo apresentada na Equação~\eqref{eq:gmm_objetivo}. No modelo sobreidentificado, a estatística $J$ de Hansen é calculada conforme as Equações~\eqref{eq:hansen_j} e~\eqref{eq:hansen_dist}.
A fim de compreender o porquê de o modelo exatamente identificado proporciona uma matriz de ponderação do GMM que não altera a estimação pontual, faz-se interessante perceber as diferenças conceituais e de resultados entre os modelos a) e b), bem como visualizar o teste Hansen feito de acordo com o modelo sobreidentificado. De modo que, ambos os modelos produzirão coeficientes e testes de validação dos instrumentos com precisões diferentes.

\begin{table}[H]
    \centering
    \small
    \caption{Resultados do GMM e teste de Hansen}
    \label{tab:gmm_hansen}
    \setlength{\tabcolsep}{5pt}
    \begin{tabular}{lcccccc}
        \hline
        Modelo 
        & $N$ 
        & $\hat{\beta}_p$ 
        & EP 
        & IC 95\% 
        & Hansen $J$ 
        & $p$-valor \\
        \hline
        GMM $Z_1$ 
        & 52 
        & $-1{,}711$ 
        & $0{,}941$ 
        & $[-3{,}555;\ 0{,}133]$ 
        & -- 
        & -- \\

        GMM $Z_2$ 
        & 52 
        & $-0{,}282$ 
        & $0{,}078$ 
        & $[-0{,}435;\ -0{,}129]$ 
        & $14{,}592$ 
        & $0{,}0001$ \\
        \hline
    \end{tabular}

    \vspace{0.3em}
    \begin{minipage}{0.92\textwidth}
    \footnotesize
    Nota: $Z_1=\{z_t\}$ e $Z_2=\{z_t,z_t^2\}$. No modelo $GMM\ Z_1$, o sistema é exatamente identificado, pois há um instrumento excluído para uma variável endógena. Por isso, não há restrições sobreidentificadoras a serem testadas e a estatística de Hansen $J$ não se aplica. No modelo $GMM\ Z_2$, há uma restrição sobreidentificadora.
    \end{minipage}
\end{table}

No caso exatamente identificado, representado pelo modelo que usa apenas o conjunto de instrumentos $Z_1= \{Z_t\}$ de numeração igual ao número de regressores estruturais endógenos, não restam momentos para serem testados, de modo que não só o teste J de Hansen resulta em valores NA - dado que graus de liberdade iguais a 0 - como também esse sistema é de solução única, não fazendo diferença optar pela escolha da matriz de ponderação já que não altera a estimativa pontual. Em contrapartida, o modelo, no qual é sobreidentificado, apresenta maior quantia de instrumentos explicando a variável endógena a serem analisados e, por isso, abre possibilidade de testar as restrições sobreidentificadoras.

Nesse sentido, em observância aos resultados adquiridos no \textit{output}, o parâmetro de elasticidade-preço evolui consideravelmente a partir da adição de um instrumento em potencial. Isso, porque a estimativa de beta em a) era de -2,93 com um erro-padrão de 2,11 - considerado elevado em demasia -, levando a um p-valor de 0,166 condizente a um coeficiente não significante estatisticamente. Aprofundando e olhando para o intervalo de confiança, percebe-se resultados alongados entre -7,06 e 1,21, o que reitera a possibilidade de ausência de efeito. Por outro lado, a estimação do modelo sobreidentificado retornou um coeficiente estimado, aproximadamente, em -1,39 acompanhado de um erro-padrão menor de 0,134. Dessa forma, o p-valor, de 3.02 x 10-25, aponta como sendo um estimador significante.

Apesar dessas observações, a validade dos instrumentos parece indicar um caminho contrário ao da ortogonalidade desejada. Isso, pelo fato de o modelo em b) apresentar estimativas consideravelmente precisas, mas com testes em sua maneira dúbios. No teste J de Hansen para o GMM de Z2, o resultado é uma estatística de 6,557 conjuntamente a um p-valor de 0,0104, o que, pela definição, ser menor que os 0,05 rejeita a hipótese nula do conjunto de instrumentos ser inteiramente exógeno. Por esse motivo, deve-se levar em conta que ao menos um dos instrumentos utilizados para análise é endógeno e, portanto, está correlacionado com o termo de erro, além de não ser possível afirmar em completude a adequação das especificações ao modelo.

\item \textbf{Instrumentos fracos.} Rode o teste $F$ efetivo de Montiel Olea e
    Pflueger para a especificação com $z_t$ como instrumento excluído. Informe os
    valores críticos utilizados pelo software ou pela tabela de referência e conclua
    se há evidência de instrumento fraco. Compare esse diagnóstico com o $F$ usual
    do primeiro estágio.
    
\textbf{Resultados}

A especificação analisada usa o primeiro estágio da Equação~\eqref{eq:primeiro_estagio}. O teste $F$ usual é o da Equação~\eqref{eq:f_usual}, enquanto o diagnóstico robusto de Montiel Olea e Pflueger segue a estatística $F^{Ef}$ e a regra de decisão das Equações~\eqref{eq:f_efetivo} e~\eqref{eq:regra_mop}. A instabilidade gerada por instrumentos fracos também pode ser entendida pela razão IV da Equação~\eqref{eq:iv_ratio}.

    O Teste F Efetivo ajusta a estatística F do primeiro estágio, que só é efetiva em um caso sem heterocedasticidade e sem autocorrelação. Basicamente, se ele for maior do que o valor crítico, rejeitamos a hipótese nula de instrumentos fracos.


\text{1. Resultados da Estimação para $Z_1 = \{z_t\}$}

A Tabela~\ref{tab:mop_z1} sintetiza as estatísticas do primeiro estágio e os valores críticos associados ao teste de Montiel Olea e Pflueger (MOP).

\begin{table}[H]
    \centering
    \caption{Resultados do Teste de Instrumentos Fracos para $Z_1$}
    \label{tab:mop_z1}
    \begin{tabular}{lc}
        \toprule
        \textbf{Estatística / Valor Crítico} & \textbf{Valor} \\
        \midrule
        Número de Observações ($N$) & 52 \\
        $F$ usual do 1º estágio & \textbf{2,661} \quad ($p = 0,109$) \\
        $F$ efetivo de Montiel Olea \& Pflueger (MOP) & \textbf{2,661} \\
        \midrule
        Valor Crítico MOP 5\% ($\tau = 5\%$) & \textbf{37,418} \\
        Valor Crítico MOP 10\% ($\tau = 10\%$) & \textbf{23,109} \\
        Valor Crítico MOP 20\% ($\tau = 20\%$) & \textbf{15,062} \\
        \bottomrule
    \end{tabular}
\end{table}

\text{2. Equivalência Algébrica: $F_{\text{usual}} \approx F_{\text{eff}}$ para $k = 1$}
\begin{figure}[H]
    \centering
    \includegraphics[width=0.88\textwidth]{Lista_1/Figures/Stata/stata_fig10_iv_curve_actual_data.pdf}
    \caption{Efeito do enfraquecimento do primeiro estágio sobre a elasticidade-preço estimada por IV. A curva foi construída a partir dos dados da demanda por frango, mantendo fixa a forma reduzida estimada, $\hat{\delta}$, e variando contrafactualmente o coeficiente do primeiro estágio, $\gamma$. Quando $\gamma$ se aproxima de zero, o denominador do estimador IV torna-se pequeno e a elasticidade-preço implícita torna-se instável.}
    \label{fig:iv_curve_actual_data}
\end{figure}
Quando há \textbf{apenas um instrumento} ($k = 1$), a estatística $F$ efetiva de MOP coincide numericamente com a estatística $F$ usual robusta do primeiro estágio. 
A divergência prática e conceitual entre os dois testes manifesta-se de forma relevante quando existem \textbf{múltiplos instrumentos} ($k \geq 2$). Sob condições de heterocedasticidade, o $F$ usual agrega a força dos instrumentos utilizando uma ponderação que ignora as correlações não-constantes dos erros, o que pode inflar artificialmente a estatística e gerar diagnósticos falsamente otimistas (falsos positivos de instrumentos fortes). O $F$ efetivo corrige essa distorção.

\text{3. Análise dos Valores Críticos e Conclusão}

Os valores críticos propostos por Montiel Olea \& Pflueger (2013) estabelecem limites rígidos de tolerância baseados no viés relativo máximo do estimador de Mínimos Quadrados em 2 Estágios (2SLS) em relação ao viés do estimador de Mínimos Quadrados Ordinários (OLS):
\begin{itemize}
    \item \textbf{$\text{CV}_{5\%} = 37,42$}: Garante um viés relativo máximo do 2SLS de apenas 5\% em relação ao OLS (ao nível de 5\% de significância).
    \item \textbf{$\text{CV}_{10\%} = 23,11$}: Garante um viés relativo máximo de 10\%.
    \item \textbf{$\text{CV}_{20\%} = 15,06$}: Garante um viés relativo máximo de 20\%.
\end{itemize}

Estes valores são substancialmente superiores ao limite tradicional de $F > 10$.

\textbf{Conclusão do Teste:} Como a estatística $F$ efetiva calculada (\textbf{2,661}) é drasticamente inferior até mesmo ao valor crítico mais tolerante ($\text{CV}_{20\%} = 15,062$), \textbf{não se rejeita a hipótese nula de instrumentos fracos}. Conclui-se que o preço real do milho contemporâneo ($z_t$) é um instrumento fraco. Consequentemente, as estimativas de elasticidade via 2SLS sob esta especificação sofrem de severo viés e erro padrão expandido, tornando a inferência convencional não confiável para fins de análise econômica ou antitruste.

A princípio, o preço do milho parece ser um bom candidato: se o milho sobe, a ração fica mais cara, o produtor do frango gasta mais e repassa esse custo para o consumidor final.
No entanto, o teste nos diz que, nesta amostra anual, essa correlação não foi forte o suficiente. Algumas razões econômicas para isso podem ser:
\\\textbf{1. Defasagem temporal:} produtores podem não repassar custos automaticamente, seja por contratos longos, ou para não perderem fatia de mercado.
\\\textbf{2. Diluição de custos:} o milho é apenas um custos. Outros fatores como energia, logística e outros grãos podem estar abafando o impacto do milho no preço final.

\item \textbf{Heterocedasticidade e teste F usual.} Discuta por que o teste $F$
    usual do primeiro estágio, derivado sob homocedasticidade, pode ser inadequado
    neste problema. Existe motivação econômica ou empírica para suspeitar de
    heterocedasticidade em uma demanda agregada anual com renda, preços e consumo
    variando ao longo do tempo? Justifique.
\textbf{Resultados}

A discussão se refere ao uso do teste $F$ usual do primeiro estágio, apresentado na Equação~\eqref{eq:f_usual}, aplicado ao primeiro estágio da Equação~\eqref{eq:primeiro_estagio}. Quando há heterocedasticidade, a comparação com a regra robusta de Montiel Olea e Pflueger, indicada na Equação~\eqref{eq:regra_mop}, torna-se mais adequada.

O teste $F$ usual do primeiro estágio é tradicionalmente derivado sob a hipótese de homocedasticidade, isto é, sob a suposição de que a variância do erro do primeiro estágio é constante ao longo das observações. Quando essa hipótese não é plausível, a estatística $F$ convencional pode deixar de refletir adequadamente a força efetiva dos instrumentos. Em particular, na presença de heterocedasticidade, o teste pode sugerir uma força do instrumento maior do que a real, produzindo uma falsa sensação de segurança sobre a identificação do modelo.

Esse problema é especialmente relevante em modelos de variáveis instrumentais, pois a inferência por 2SLS é muito sensível à força do primeiro estágio. Um instrumento pode parecer estatisticamente relevante pelo teste $F$ usual, mas ainda assim ser fraco quando a variância dos erros não é constante. Por isso, a literatura recomenda cautela com a regra de bolso $F>10$, principalmente quando há heterocedasticidade. Nesses casos, estatísticas robustas, como o $F$ efetivo de Montiel Olea e Pflueger, são mais adequadas para diagnosticar instrumentos fracos.

Do ponto de vista econômico, há forte motivação para suspeitar de heterocedasticidade em uma demanda agregada anual. A amostra cobre um período longo, no qual renda, preços relativos e consumo variam de maneira expressiva. Ao longo do tempo, choques macroeconômicos, mudanças de regime inflacionário, alterações tecnológicas, transformações no mercado de alimentos e mudanças nos padrões de consumo podem afetar não apenas o nível médio da demanda, mas também a variância dos choques não observados. Assim, não é razoável supor, sem verificação, que a dispersão dos erros seja constante durante todo o período.

Essa suspeita também aparece visualmente na Figura~\ref{fig:residuos_mqo_lowess}, que apresenta os resíduos do MQO contra os valores ajustados. Se a especificação linear em logaritmos estivesse capturando bem a relação média e se a variância dos erros fosse aproximadamente constante, esperaríamos resíduos distribuídos de forma aleatória ao redor de zero, sem padrão sistemático. No entanto, a curva LOWESS mostra um formato não linear: os resíduos ficam próximos de zero em parte da amostra, tornam-se negativos em uma região intermediária dos valores ajustados e depois voltam a crescer nos valores ajustados mais altos. Esse padrão sugere que pode haver não linearidade, mudança na dispersão dos resíduos ou alguma forma de má especificação funcional.

\begin{figure}[H]
    \centering
    \includegraphics[width=0.88\textwidth]{Lista_1/Figures/Stata/stata_fig03_residuals_fitted.pdf}
    \caption{Resíduos do MQO contra valores ajustados, com curva não linear LOWESS. A linha horizontal indica resíduo igual a zero. A curva suavizada auxilia na verificação visual de padrões sistemáticos nos resíduos, como indícios de não linearidade, heterocedasticidade ou má especificação funcional.}
    \label{fig:residuos_mqo_lowess}
\end{figure}

Portanto, há tanto uma motivação teórica quanto uma evidência empírica preliminar para tratar a homocedasticidade com cautela neste problema. A presença de padrões nos resíduos reforça a decisão de usar erros-padrão robustos à heterocedasticidade e de complementar o diagnóstico do primeiro estágio com medidas robustas, como o $F$ efetivo de Montiel Olea e Pflueger. Dessa forma, a análise evita depender exclusivamente do teste $F$ usual, que pode ser inadequado quando os erros apresentam variância não constante.

\item \textbf{Especificações alternativas de instrumentos.} Estime o primeiro
    estágio e a equação estrutural por 2SLS para os seguintes conjuntos de
    instrumentos excluídos:
    \[
    \begin{aligned}
        Z_1 &= \{z_t\},\\
        Z_2 &= \{z_t, z_t^2\},\\
        Z_3 &= \{z_t, z_t^2, z_t^3\},\\
        Z_4 &= \{z_{t-1}\},\\
        Z_5 &= \{z_{t-1}, z_{t-1}^2\},\\
        Z_6 &= \{z_t, z_{t-1}\},\\
        Z_7 &= \{z_t, z_t^2, z_{t-1}, z_{t-1}^2\}.
    \end{aligned}
    \]
    Como $z_t \equiv \ln(P COR_t/CPI_t)$, o termo quadrático em log é
    $z_t^2 = [\ln(P COR_t/CPI_t)]^2$.
\textbf{Resultados}

As especificações estimadas correspondem aos conjuntos de instrumentos definidos na Equação~\eqref{eq:conjuntos_instrumentos}. Para cada conjunto, estima-se o primeiro estágio geral da Equação~\eqref{eq:base_primeiro_estagio} e a equação estrutural de demanda da Equação~\eqref{eq:base_estrutural}, usando o estimador 2SLS da Equação~\eqref{eq:2sls}.

Para avaliar as especificações alternativas, foram estimados o primeiro estágio e a equação estrutural por 2SLS para os sete conjuntos de instrumentos propostos. No primeiro estágio, a variável endógena é o preço real do frango, $p^{ch}_t$, e os instrumentos excluídos variam entre o preço real contemporâneo do milho, suas potências, defasagens e combinações entre termos contemporâneos e defasados.

A Tabela~\ref{tab:q8_instrumentos_alt} resume os principais resultados. As colunas de $R^2$ parcial e $F$ usual dizem respeito à força dos instrumentos no primeiro estágio. Já a coluna $\hat{\beta}_p$ apresenta a elasticidade-preço da demanda estimada por 2SLS na equação estrutural.

\begin{table}[H]
    \centering
    \caption{Especificações alternativas de instrumentos: primeiro estágio e 2SLS}
    \label{tab:q8_instrumentos_alt}
    \begin{tabular}{c c c c c c}
        \hline
        Instrumentos & $N$ & $R^2$ parcial & $F$ usual & $\hat{\beta}_p$ & IC 95\% \\
        \hline
        $Z_1=\{z_t\}$ 
        & 52 & 0{,}041 & 2{,}661 & -1{,}711 & [-3{,}555;\ 0{,}133] \\

        $Z_2=\{z_t,z_t^2\}$ 
        & 52 & 0{,}605 & 29{,}374 & -0{,}261 & [-0{,}415;\ -0{,}108] \\

        $Z_3=\{z_t,z_t^2,z_t^3\}$ 
        & 52 & 0{,}627 & 23{,}224 & -0{,}256 & [-0{,}401;\ -0{,}111] \\

        $Z_4=\{z_{t-1}\}$ 
        & 51 & 0{,}018 & 0{,}998 & -2{,}739 & [-7{,}519;\ 2{,}041] \\

        $Z_5=\{z_{t-1},z_{t-1}^2\}$ 
        & 51 & 0{,}498 & 24{,}222 & -0{,}275 & [-0{,}478;\ -0{,}071] \\

        $Z_6=\{z_t,z_{t-1}\}$ 
        & 51 & 0{,}062 & 1{,}724 & -1{,}233 & [-2{,}460;\ -0{,}006] \\

        $Z_7=\{z_t,z_t^2,z_{t-1},z_{t-1}^2\}$ 
        & 51 & 0{,}610 & 17{,}309 & -0{,}269 & [-0{,}435;\ -0{,}103] \\
        \hline
    \end{tabular}
\end{table}

Os resultados mostram diferenças importantes entre os conjuntos de instrumentos. A especificação $Z_1$, que utiliza apenas o preço real contemporâneo do milho, apresenta baixo poder explicativo no primeiro estágio: o $R^2$ parcial é de apenas 0{,}041 e o $F$ usual é igual a 2{,}661. A estimativa estrutural obtida com esse instrumento é $\hat{\beta}_p=-1{,}711$, com intervalo de confiança amplo e que inclui zero. Isso sugere que, embora o sinal seja economicamente esperado, a estimativa é imprecisa e pode estar sendo afetada por instrumento fraco.

A inclusão do termo quadrático em $Z_2=\{z_t,z_t^2\}$ altera fortemente o primeiro estágio. O $R^2$ parcial sobe para 0{,}605 e o $F$ usual passa para 29{,}374. A elasticidade estimada também muda substancialmente, tornando-se $\hat{\beta}_p=-0{,}261$, com intervalo de confiança inteiramente negativo. Esse resultado indica uma demanda inelástica e muito próxima da estimativa obtida por MQO, mas agora com instrumentos aparentemente mais relevantes no primeiro estágio.

A especificação $Z_3$, que acrescenta o termo cúbico $z_t^3$, mantém um primeiro estágio forte, com $R^2$ parcial de 0{,}627 e $F$ usual de 23{,}224. A elasticidade estimada é $\hat{\beta}_p=-0{,}256$, praticamente igual à obtida com $Z_2$. Isso sugere que a inclusão do termo cúbico não altera de forma substantiva a estimativa estrutural em relação à especificação com termo quadrático.

No caso dos instrumentos defasados, $Z_4=\{z_{t-1}\}$ apresenta desempenho fraco. O $R^2$ parcial é de apenas 0{,}018 e o $F$ usual é 0{,}998. A estimativa da elasticidade é $\hat{\beta}_p=-2{,}739$, mas o intervalo de confiança é extremamente amplo, indo de $-7{,}519$ a $2{,}041$. Esse resultado indica forte imprecisão e reforça a evidência de que o preço real defasado do milho, isoladamente, é um instrumento fraco para o preço real do frango.

A especificação $Z_5=\{z_{t-1},z_{t-1}^2\}$ melhora muito o primeiro estágio em relação a $Z_4$. O $R^2$ parcial sobe para 0{,}498 e o $F$ usual para 24{,}222. A elasticidade estimada é $\hat{\beta}_p=-0{,}275$, com intervalo de confiança negativo. Assim como em $Z_2$ e $Z_3$, a adição de um termo não linear aumenta bastante a relevância estatística dos instrumentos e produz uma estimativa de demanda inelástica.

A especificação $Z_6=\{z_t,z_{t-1}\}$ apresenta resultado intermediário, mas ainda problemático. Apesar de combinar o preço contemporâneo e defasado do milho, o $R^2$ parcial é de apenas 0{,}062 e o $F$ usual é 1{,}724. A elasticidade estimada é $\hat{\beta}_p=-1{,}233$, com intervalo de confiança bastante amplo. Isso indica que a simples combinação linear entre o instrumento contemporâneo e o defasado não é suficiente para gerar um primeiro estágio forte.

Por fim, a especificação $Z_7=\{z_t,z_t^2,z_{t-1},z_{t-1}^2\}$ apresenta um primeiro estágio relativamente forte, com $R^2$ parcial de 0{,}610 e $F$ usual de 17{,}309. A elasticidade estimada é $\hat{\beta}_p=-0{,}269$, novamente próxima das estimativas obtidas com $Z_2$, $Z_3$ e $Z_5$. Esse resultado sugere que os conjuntos de instrumentos com termos quadráticos geram estimativas mais estáveis e concentradas em torno de uma elasticidade-preço inelástica.

Em conjunto, os resultados mostram que as especificações com apenas um instrumento linear, especialmente $Z_1$ e $Z_4$, apresentam sinais de fraqueza no primeiro estágio e produzem estimativas mais instáveis. Já as especificações que incluem termos quadráticos, como $Z_2$, $Z_3$, $Z_5$ e $Z_7$, aumentam substancialmente a relevância dos instrumentos e levam a estimativas de $\beta_p$ próximas de $-0{,}26$ a $-0{,}28$. No entanto, esse ganho de relevância estatística deve ser interpretado com cautela: a inclusão de termos não lineares pode melhorar o ajuste do primeiro estágio, mas também torna a defesa econômica da hipótese de exclusão mais delicada, pois é necessário sustentar que não apenas $z_t$, mas também suas transformações e defasagens, afetam a demanda por frango apenas por meio do preço do frango.

\item \textbf{Tabela comparativa.} Construa uma tabela com os resultados dos
    modelos do item anterior. A tabela deve conter, no mínimo, as seguintes colunas:
    \[
    \begin{array}{c|c|c|c|c|c|c}
    \text{Instrumentos} 
    & N 
    & \hat{\beta}_p 
    & \text{IC 95\%} 
    & F\text{ usual} 
    & F\text{ efetivo MOP} 
    & \text{Hansen } J \\
    \hline
    Z_1 
    & 52 
    & -1{,}711 
    & [-3{,}555;\ 0{,}133] 
    & 2{,}661 
    & 2{,}661 
    & \text{--} \\

    Z_2 
    & 52 
    & -0{,}261 
    & [-0{,}415;\ -0{,}108] 
    & 29{,}374 
    & 28{,}914 
    & 14{,}592 \\

    Z_3 
    & 52 
    & -0{,}256 
    & [-0{,}401;\ -0{,}111] 
    & 23{,}224 
    & 23{,}975 
    & 14{,}683 \\

    Z_4 
    & 51 
    & -2{,}739 
    & [-7{,}519;\ 2{,}041] 
    & 0{,}998 
    & 0{,}998 
    & \text{--} \\

    Z_5 
    & 51 
    & -0{,}275 
    & [-0{,}478;\ -0{,}071] 
    & 24{,}222 
    & 22{,}863 
    & 14{,}519 \\

    Z_6 
    & 51 
    & -1{,}233 
    & [-2{,}460;\ -0{,}006] 
    & 1{,}724 
    & 1{,}547 
    & 1{,}996 \\

    Z_7 
    & 51 
    & -0{,}269 
    & [-0{,}435;\ -0{,}103] 
    & 17{,}309 
    & 12{,}484 
    & 15{,}699 \\
\end{array}
\]
\noindent
Nota: nos modelos exatamente identificados, como $Z_1$ e $Z_4$, o teste de Hansen $J$ não se aplica, pois não há restrições sobreidentificadoras a serem testadas.
   Com base nos resultados dos modelos definidos na Equação~\eqref{eq:conjuntos_instrumentos}, interprete como as estimativas de $\beta_p$ mudam quando os instrumentos se
    tornam mais ou menos fortes. Discuta também se a adição de termos não lineares
    aumenta relevância, mas pode tornar a hipótese de exclusão menos plausível ou
    apenas mais difícil de defender.

A comparação entre as colunas da tabela deve ser lida à luz do teste $F$ usual da Equação~\eqref{eq:f_usual}, do $F$ efetivo de MOP da Equação~\eqref{eq:f_efetivo} e do teste $J$ de Hansen das Equações~\eqref{eq:hansen_j} e~\eqref{eq:hansen_dist}.

\item \textbf{Teste J de Hansen e validade dos instrumentos.} Para os modelos
    sobreidentificados, acrescente o valor-$p$ do teste $J$ de Hansen. Explique
    cuidadosamente a diferença entre: (i) instrumento relevante; (ii) instrumento
    exógeno; (iii) instrumento forte; e (iv) conjunto de instrumentos válido. O teste
    $J$ avalia a validade conjunta das restrições sobreidentificadoras, mas não prova
    que todos os instrumentos são exógenos.
\textbf{Resultados}

A distinção conceitual abaixo remete às condições de relevância e exogeneidade das Equações~\eqref{eq:relevancia} e~\eqref{eq:exogeneidade}. Para os modelos sobreidentificados, o teste $J$ segue as Equações~\eqref{eq:hansen_j} e~\eqref{eq:hansen_dist}, tendo como hipótese nula a condição de momento indicada na Equação~\eqref{eq:h0_hansen}.

Diferenciando os conceitos:

\begin{enumerate}[label=(\roman*)]
    \item \textbf{Instrumento relevante:} 
    O instrumento será considerado relevante se tiver correlação com a variável endógena, ou seja, se 
    \[
        \operatorname{Corr}(X_i,Z_i)\neq 0,
    \]
    sendo $X_i$ a variável endógena e $Z_i$ o instrumento.\footnote{GOMES, Victor. \textit{Introdução à Economia Industrial e ao Antitruste}. Notas de aula, Universidade de Brasília.} 
    No contexto da lista de exercícios, os instrumentos são relevantes se suas covariâncias individuais com o preço real do frango, $p_t^{ch}$, não são nulas.

    \item \textbf{Instrumento exógeno:} 
    A correlação entre o instrumento e o termo de erro deve ser nula:
    \[
        \operatorname{Corr}(Z_i,\varepsilon_i)=0.
    \]
    \footnote{GOMES, Victor. \textit{Introdução à Economia Industrial e ao Antitruste}. Notas de aula, Universidade de Brasília.}
    Portanto, os instrumentos utilizados, como, por exemplo, o preço do milho ($z_t$), devem ser não correlacionados ao termo de erro. Em outras palavras, sua variação deve ser exógena ao modelo utilizado.

    As duas condições mencionadas acima configuram a validade dos instrumentos: relevância e exogeneidade.

    \item \textbf{Instrumento forte:} 
    Para averiguar se o instrumento é forte, pode-se usar a regra de bolso proposta por Stock e Yogo (2005)\footnote{Stock, James H. e Motohiro Yogo. “Testing for Weak Instruments in Linear IV Regression.” In: Donald W. Andrews (org.), \textit{Identification and Inference for Econometric Models}. New York, Cambridge University Press, 2005.} 
    e Staiger e Stock (1997)\footnote{Staiger, Douglas e James H. Stock. “Instrumental Variables Regression with Weak Instruments.” \textit{Econometrica}, 65 (3), 1997.}, que indica que o(s) instrumento(s) será(ão) considerado(s) forte(s) se a estatística $F>10$.\footnote{GOMES, Victor. \textit{Introdução à Economia Industrial e ao Antitruste}. Notas de aula, Universidade de Brasília.} 
    Contudo, essa regra de bolso só pode ser utilizada em cenários de uma única variável endógena e erros de primeiro estágio homocedásticos. No caso de erro heterocedástico, é possível utilizar a estatística $F^{Ef}$, como foi solicitado em outros exercícios da lista.

    \item \textbf{Conjunto de instrumentos válido:} 
    Como mencionado acima, um instrumento é considerado válido se cumpre concomitantemente as condições de relevância e exogeneidade. Todavia, em modelos sobreidentificados, é necessário testar a validade conjunta dos instrumentos, o que é justamente feito pelo teste $J$ de Hansen. É relevante a ressalva feita pelo próprio enunciado: esse teste não prova a validade individual dos instrumentos, mas avalia a validade conjunta das restrições sobreidentificadoras.
\end{enumerate}

A Tabela~\ref{tab:hansen_sobreidentificados} apresenta os resultados do teste $J$ de Hansen para os modelos sobreidentificados. Os modelos $Z_1$ e $Z_4$ não aparecem na tabela porque são exatamente identificados, isto é, possuem o mesmo número de instrumentos excluídos e variáveis endógenas; nesses casos, não há restrições sobreidentificadoras a serem testadas.

\begin{table}[H]
    \centering
    \small
    \caption{Teste $J$ de Hansen para modelos sobreidentificados}
    \label{tab:hansen_sobreidentificados}
    \setlength{\tabcolsep}{5pt}
    \begin{tabular}{lcccc}
        \hline
        Modelo 
        & Instrumentos excluídos 
        & Hansen $J$ 
        & \textit{gl} 
        & Valor-$p$ \\
        \hline
        $Z_2$ 
        & $\{z_t,\ z_t^2\}$ 
        & 14{,}592 
        & 1 
        & 0{,}0001 \\

        $Z_3$ 
        & $\{z_t,\ z_t^2,\ z_t^3\}$ 
        & 14{,}683 
        & 2 
        & 0{,}0006 \\

        $Z_5$ 
        & $\{z_{t-1},\ z_{t-1}^2\}$ 
        & 14{,}519 
        & 1 
        & 0{,}0001 \\

        $Z_6$ 
        & $\{z_t,\ z_{t-1}\}$ 
        & 1{,}996 
        & 1 
        & 0{,}1577 \\

        $Z_7$ 
        & $\{z_t,\ z_t^2,\ z_{t-1},\ z_{t-1}^2\}$ 
        & 15{,}699 
        & 3 
        & 0{,}0013 \\
        \hline
    \end{tabular}

    \vspace{0.3em}
    \begin{minipage}{0.92\textwidth}
    \footnotesize
    Nota: $gl$ indica os graus de liberdade do teste, dados pelo número de instrumentos excluídos menos o número de variáveis endógenas. Os modelos $Z_1$ e $Z_4$ são exatamente identificados e, por isso, não possuem restrições sobreidentificadoras a serem testadas.
    \end{minipage}
\end{table}

Pela Tabela~\ref{tab:hansen_sobreidentificados}, observa-se que os modelos $Z_2$, $Z_3$, $Z_5$ e $Z_7$ apresentam valores-$p$ inferiores a 5\%. Portanto, para essas especificações, rejeita-se a hipótese nula de validade conjunta das restrições sobreidentificadoras. Isso sugere que ao menos um dos instrumentos utilizados pode estar correlacionado com o termo de erro, ou que a especificação do conjunto de instrumentos não é plenamente compatível com a hipótese de exogeneidade conjunta.

O modelo $Z_6$, por sua vez, apresenta valor-$p$ igual a 0{,}1577. Como esse valor é superior a 5\%, não se rejeita a hipótese nula do teste $J$ de Hansen. Entretanto, isso não significa que seus instrumentos estejam comprovadamente corretos ou individualmente exógenos. Significa apenas que, nessa especificação, o teste não encontrou evidência estatística suficiente contra a validade conjunta das restrições sobreidentificadoras.

\item \textbf{Inferência robusta a instrumentos fracos.} Para pelo menos três
    especificações de instrumentos, reporte também um intervalo de confiança robusto
    a instrumentos fracos, como Anderson--Rubin ou CLR, se o software utilizado
    estiver disponível. Compare esses intervalos com os intervalos convencionais
    baseados na aproximação normal do 2SLS.

\textbf{Resultados}

Os intervalos convencionais partem da estimação 2SLS da Equação~\eqref{eq:2sls}. Já os intervalos robustos a instrumentos fracos seguem a lógica do teste Anderson--Rubin, descrita na Equação~\eqref{eq:ar}.

% Requires: \usepackage{amsmath}
\begin{table}[h]
    \centering
    \caption{Intervalos de Confiança 2SLS}
    \label{tab:q11_ic_2sls}
    \begin{tabular}{cccc}
        \hline
        Especificação & Instrumentos Excluídos & Estatística $F$ & 2SLS \\
        \hline
        $Z_{1}$ & $\{ z_{t} \}$ & 2,0675  & $[-3{,}8338,\; 0{,}4121]$ \\
        $Z_{2}$ & $\{ z_{t}, z_{t}^{2} \}$ & 36,0657 & $[-0{,}4374,\; -0{,}0855]$ \\
        $Z_{3}$ & $\{ z_{t}, z_{t}^{2}, z_{t}^{3} \}$ & 25,7382 & $[-0{,}4289,\; -0{,}0829]$ \\
        \hline
    \end{tabular}
\end{table}
% Requires: \usepackage{amsmath}
\begin{table}[h]
    \centering
    \caption{Intervalos de Confiança AR e CLR}
    \label{tab:q11_ic_ar_clr}
    \begin{tabular}{c c c}
        \hline
        Especificação & AR & CLR \\
        \hline
        $Z_{1}$ & $(-\infty, -0{,}7606] \cup [3{,}5134, \infty)$ & $(-\infty, -0{,}7606] \cup [3{,}5134, \infty)$ \\
        $Z_{2}$ & Conjunto Vazio & $[-0{,}4472, -0{,}0601]$ \\
        $Z_{3}$ & Conjunto Vazio & $[-0{,}4236, -0{,}0496]$ \\
        \hline
    \end{tabular}
\end{table}

    No caso da especificação \textit{Z\textsubscript{1}}, a estatística \textit{F} do primeiro estágio é muito baixa, significando um instrumento fraco. Por se tratar de um modelo exatamente identificado, os intervalos de confiança, ajustados para instrumentos fracos por meio dos testes de Anderson-Rubin (AR) e Razão de Verossimilhança Condicional (CLR), são exatamente iguais, ilimitados e disjuntos. Ao contrário, as outras especificações possuem estatísticas F elevadas, que por sua vez podem levar a interpretações incorretas. Seus intervalos estimados por AR e CLR são conjuntos fechados, mais próximos do formato pelo qual estamos tradicionalmente acostumados. 

    No que concerne aos intervalos convencionais baseados na aproximação normal do 2SLS, temos que essa falha grosseiramente ao estimar por instrumentos fracos. Os estimadores, com isso, são enviesados e os intervalos imprecisos, bem como o apresentado pela especificação \textit{Z\textsubscript{1}}. Com o propósito de solucionar esse problema, aplica-se as geometrias de AR e CLR. A cobertura de 95\% em grandes amostras é preservada por esses testes, independentemente do tamanho da amostra. 

    Como discutido anteriormente, apesar das magnitudes dispostas obedecerem a regra de bolso, as variáveis exógenas provavelmente não satisfazem a condição de exogeneidade, podendo refletir nos conjuntos vazios do teste AR. Entende-se por ele que não existem parâmetros que tornem as restrições sobre-identificadoras aceitas, rejeitando universalmente o modelo. Entretanto, se a abordagem adotada for proveniente do teste CRL, argumenta-se os resultados a partir dos intervalos fechados da tabela. Desse modo, a preferência pelos testes oferece visões distintas. O teste AR é geralmente mais preciso e robusto. Por outro lado, costuma a ser ineficiente para muitos instrumentos sobre-identificados.

\item \textbf{Interpretação de economia industrial.} A partir dos resultados,
    responda: a elasticidade-preço estimada é compatível com uma demanda
    inelástica ou elástica? Como essa conclusão afetaria uma análise de poder de
    mercado, margem ou perda crítica? Explique por que uma elasticidade estimada
    com instrumento fraco pode ser enganosa para análise antitruste.


\textbf{Resultados}

A interpretação de economia industrial conecta a elasticidade estimada na demanda da Equação~\eqref{eq:demanda_estrutural} à regra de margem da Equação~\eqref{eq:lerner} e à expressão de perda crítica da Equação~\eqref{eq:critical_loss}. A cautela com instrumentos fracos decorre da razão IV apresentada na Equação~\eqref{eq:iv_ratio}.

Após reunir os resultados das questões anteriores, a interpretação da elasticidade-preço da demanda por frango deve ser feita com cautela. A estimativa por MQO, reportada anteriormente, foi de aproximadamente $\hat{\beta}_p=-0{,}265$, o que sugere uma demanda inelástica, pois $|\hat{\beta}_p|<1$. Nessa leitura, um aumento de 1\% no preço real do frango estaria associado a uma queda de apenas 0,265\% na quantidade demandada, mantendo constantes renda e preço da carne bovina.

No entanto, como discutido nas questões sobre endogeneidade, o MQO pode estar enviesado. Em uma equação de demanda, o preço do frango pode ser determinado simultaneamente com a quantidade, de modo que choques positivos de demanda elevam tanto o preço quanto a quantidade de equilíbrio. Esse mecanismo tende a gerar viés positivo no coeficiente de preço, aproximando a elasticidade estimada de zero e fazendo a demanda parecer mais inelástica do que realmente é.

A comparação visual da Figura~\ref{fig:mqo_2sls} mostra justamente que as estimativas de $\beta_p$ variam bastante conforme o método e o conjunto de instrumentos utilizado. As especificações com instrumentos lineares fracos, como $Z_1$ e $Z_4$, produzem elasticidades mais altas em módulo, sugerindo uma demanda elástica. Contudo, esses modelos também apresentam intervalos de confiança amplos e forte evidência de instrumentos fracos. Por isso, não é adequado concluir mecanicamente que a demanda é elástica apenas com base nessas estimativas.

Por outro lado, as especificações com termos quadráticos, como $Z_2$, $Z_3$, $Z_5$ e $Z_7$, retornam elasticidades próximas de $-0{,}26$ a $-0{,}28$, compatíveis com uma demanda inelástica. Essas estimativas são mais precisas, mas também devem ser interpretadas com cautela, pois algumas delas apresentam rejeição no teste $J$ de Hansen, indicando possível problema na validade conjunta dos instrumentos.

Assim, a conclusão mais prudente é que os resultados mais estáveis apontam para uma demanda inelástica, mas a classificação da demanda depende fortemente da especificação instrumental adotada. As estimativas que sugerem demanda elástica aparecem justamente em modelos nos quais há maior preocupação com instrumentos fracos e intervalos de confiança muito amplos.

Do ponto de vista de Economia Industrial, essa distinção é central. Se a demanda for inelástica, as firmas teriam maior capacidade de elevar preços sem perder grande parcela da quantidade vendida. Isso poderia ser interpretado como indício de maior poder de mercado, pois os consumidores reagiriam pouco a aumentos de preço. Nesse cenário, margens preço-custo mais elevadas seriam mais sustentáveis, e um aumento unilateral de preços teria maior chance de ser lucrativo.

Se, ao contrário, a demanda for elástica, aumentos de preço provocariam quedas proporcionalmente maiores na quantidade demandada. Isso limitaria o poder de mercado unilateral das firmas, pois consumidores poderiam reduzir o consumo de frango ou migrar para substitutos próximos, como outras carnes. Nesse caso, a margem sustentável seria menor, e a perda crítica associada a um aumento de preço seria mais facilmente ultrapassada.

A Figura~\ref{fig:iv_curve_actual_data} ajuda a entender por que uma elasticidade estimada com instrumento fraco pode ser enganosa para análise antitruste. O estimador IV pode ser interpretado como uma razão entre a forma reduzida e o primeiro estágio. Quando o primeiro estágio se aproxima de zero, o denominador da razão fica pequeno, tornando a elasticidade-preço estimada extremamente instável. Assim, pequenas variações nos dados ou na força do instrumento podem gerar grandes oscilações em $\hat{\beta}_p$.

Por isso, uma elasticidade estimada com instrumento fraco pode levar a conclusões equivocadas. Se o instrumento fraco gerar uma elasticidade artificialmente baixa em módulo, a autoridade antitruste pode superestimar o poder de mercado das firmas. Por outro lado, se o instrumento fraco gerar uma elasticidade excessivamente alta em módulo, pode-se subestimar esse poder de mercado. Em ambos os casos, decisões sobre margens, perda crítica, fusões ou condutas podem se tornar pouco confiáveis.

Portanto, a elasticidade-preço estimada deve ser interpretada em conjunto com os diagnósticos de força e validade dos instrumentos. Para fins de análise antitruste, não basta observar se $|\hat{\beta}_p|$ é maior ou menor que 1; é necessário avaliar se a estimativa está bem identificada, se os instrumentos são suficientemente fortes e se a hipótese de exclusão é economicamente plausível.

\item \textbf{Preço da carne bovina como variável potencialmente endógena.}
    Discuta se $p_t^b$ deve ser tratado como exógeno. Se ele for endógeno, que tipo
    de instrumento seria necessário? Avalie se exportações de carnes, $EATEX$,
    poderiam deslocar oferta de carnes sem afetar diretamente a demanda por frango.
    Deixe explícita a hipótese de exclusão.

\textbf{Resultados}

Do ponto de vista econométrico, o preço da carne bovina ($p_t^b$) não deve ser tratado como exógeno. Isto porque frango e carne bovina são bens substitutos ou, de outro modo, bens cujos equilíbrios de mercado são determinados simultaneamente. Dessa forma, qualquer choque não observado, capturado pelo termo de erro, que aumente a demanda agregada por proteínas, exercerá influência sobre ambas as curvas. Nesse sentido, não é válida a proposição $Cov(p_t^b , u_t) \neq 0$). Traduz-se dessa afirmativa a confirmação do viés de simultaneidade e, não obstante, a classificação da variável como endógena. Ao admitirmos $p_t^{ch}$ e $p_t^b$ como variáveis endógenas no modelo, a condição de ordem da identificação requer introduzir pelo menos mais um instrumento excluído. Conceitualmente, esse instrumento precisa deslocar somente a curva de oferta da carne bovina. No caso de os efeitos repousarem sobre o preço do frango, a desfaz-se a prerrogativa da inclusão dessa \textit{proxy}. 

    Em relação às exportações de carne, \textit{EATEX}, enquanto instrumento, utilizamos a análise econômica voltada às hipóteses de relevância e exogeneidade. A respeito da primeira, é razoável presumir que a dinâmica exportadora da carne bovina influi diretamente na distribuição do produto no setor interno. Pela Figura ~\ref{fig:placeholder}, temos a representação visual da interação das duas variáveis. Infere-se, portanto, que \textit{EATEX} é uma \textit{proxy} útil para captarmos choques positivos de oferta, isto é, de superprodução, na qual o excedente gerado é escoado por meio das exportações. Pelo prisma da exogeneidade, podemos nos assentar na concepção do comportamento do consumidor doméstico. Espera-se que suas preferências não sejam passíveis de mudanças significantes via dinâmica do comércio internacional e superprodução. Desse raciocínio, se depreende a não correlação do instrumento selecionado com o termo de erro. 

    Por fim, discute-se a hipótese de exclusão explícita. Essa pode ser definida pelo seguinte enunciado: As exportações de carne atuam sobre a quantidade demanda de frango no marcado doméstico estritamente através do seu peso sobre o preço de equilíbrio da carne bovina. Com isso, rejeitamos a conjectura de choques relacionados a pauta externa correlacionados com choques não observados na preferência dos consumidores por frango. 

\begin{figure}
    \centering
    \includegraphics[width=1\linewidth]{Lista_1/Figures/Stata/fig_q13_eatex_pb_first_stage.pdf}
    \caption{Relação Empírica entre Exportações de Carne e o Preço Doméstico da Carne Bovina (Primeiro Estágio) }
    \label{fig:placeholder}
\end{figure}
% A resposta substantiva da questão 13 não constava no trecho original enviado. A organização foi preservada para manter a numeração da lista.




\item \textbf{Simulação de instrumentos fracos.} Crie uma simulação simples em que a correlação entre instrumento e variável endógena seja gradualmente reduzida. Mostre como o viés de 2SLS, a dispersão do estimador e a cobertura dos intervalos de confiança convencionais se deterioram quando o instrumento fica fraco.

\textbf{Resultados}

A simulação operacionaliza o desenho das Equações~\eqref{eq:sim_primeiro_estagio} e~\eqref{eq:sim_estrutural}. A deterioração do desempenho do 2SLS é interpretada à luz da razão IV da Equação~\eqref{eq:iv_ratio} e do teste $F$ usual do primeiro estágio em~\eqref{eq:f_usual}, todo o seu processo operacional esta descrito por ~\ref{alg:simulacao_instrumentos_fracos}

A simulação foi construída para ilustrar, em ambiente controlado, o problema de instrumentos fracos. O desenho básico considera uma variável endógena explicada por um instrumento cuja força é governada pelo parâmetro $\pi$. Assim, à medida que $\pi$ diminui, a correlação entre o instrumento e a variável endógena também se reduz. O parâmetro verdadeiro usado como referência na simulação é $\beta_p=-1$.

\begin{algorithm}[H]
\DontPrintSemicolon
\caption{Simulação de instrumentos fracos}
\label{alg:simulacao_instrumentos_fracos}

\Entrada{
Conjunto de forças do instrumento 
$\Pi=\{1,\ 0{,}50,\ 0{,}25,\ 0{,}10,\ 0{,}05,\ 0{,}02\}$; 
número de repetições $R=250$; tamanho amostral $n=200$; 
parâmetro verdadeiro $\beta_p=-1$; correlação entre erros $\rho=0{,}6$.
}

\Saida{
Para cada $\pi \in \Pi$: média do 2SLS, viés do 2SLS, desvio-padrão do 2SLS, 
$F$ médio do primeiro estágio e cobertura do intervalo convencional de 95\%.
}

\Inicio{
    criar uma tabela vazia de resultados agregados\;

    \ParaCada{$\pi \in \Pi$}{
        criar vetores vazios para armazenar $\hat{\beta}^{2SLS}$, 
        $\hat{\beta}^{MQO}$, $F$ do primeiro estágio e indicador de cobertura\;

        \Para{$r=1,\ldots,R$}{
            gerar o instrumento exógeno:
            \[
                z_i \sim N(0,1), \qquad i=1,\ldots,n;
            \]

            gerar choques independentes:
            \[
                e_{1i}\sim N(0,1), 
                \qquad
                e_{2i}\sim N(0,1);
            \]

            construir os erros correlacionados:
            \[
                u_i=e_{1i},
                \qquad
                v_i=\rho e_{1i}+\sqrt{1-\rho^2}\,e_{2i};
            \]

            construir a variável endógena, conforme a Equação~\eqref{eq:sim_primeiro_estagio}:
            \[
                x_i=\pi z_i+v_i;
            \]

            construir a variável dependente, conforme a Equação~\eqref{eq:sim_estrutural}:
            \[
                y_i=\beta_p x_i+u_i;
            \]

            estimar MQO de $y_i$ em $x_i$ e armazenar $\hat{\beta}^{MQO}$\;

            estimar o primeiro estágio de $x_i$ em $z_i$ e armazenar a estatística $F$\;

            estimar 2SLS usando $z_i$ como instrumento para $x_i$ e armazenar $\hat{\beta}^{2SLS}$\;

            construir o intervalo convencional de 95\%:
            \[
                IC_{95\%}
                =
                \hat{\beta}^{2SLS}
                \pm
                1{,}96\cdot EP(\hat{\beta}^{2SLS});
            \]

            registrar se $IC_{95\%}$ contém o valor verdadeiro $\beta_p=-1$\;
        }

        calcular:
        \[
            \overline{\hat{\beta}}^{2SLS}, \qquad
            \overline{\hat{\beta}}^{MQO}, \qquad
            \text{viés}=\overline{\hat{\beta}}^{2SLS}-\beta_p,
        \]
        \[
            DP(\hat{\beta}^{2SLS}), \qquad
            \overline{F}, \qquad
            \text{cobertura}_{95\%};
        \]

        adicionar os resultados de $\pi$ à tabela agregada\;
    }

    salvar a tabela da simulação e gerar os gráficos de viés, dispersão, força do primeiro estágio e cobertura\;

    \Retorna{tabela e figuras da simulação}
}
\end{algorithm}

\begin{table}[H]
    \centering
    \caption{Simulação de instrumentos fracos}
    \label{tab:simulacao_instrumentos_fracos}
    \begin{tabular}{ccccccc}
        \hline
        $\pi$ & Repetições & Média 2SLS & Viés 2SLS & DP 2SLS & $F$ médio & Cobertura 95\% \\
        \hline
        1,00 & 250 & -0,998 &  0,002 &  0,064 & 208,871 & 96,8\% \\
        0,50 & 250 & -1,014 & -0,014 &  0,158 &  53,100 & 94,4\% \\
        0,25 & 250 & -1,100 & -0,100 &  0,392 &  13,402 & 95,6\% \\
        0,10 & 250 & -0,479 &  0,521 &  5,418 &   3,352 & 94,8\% \\
        0,05 & 250 & -4,596 & -3,596 & 58,596 &   1,662 & 92,4\% \\
        0,02 & 250 &  1,130 &  2,130 & 48,751 &   1,111 & 97,2\% \\
        \hline
    \end{tabular}
\end{table}

A Tabela~\ref{tab:simulacao_instrumentos_fracos} resume o resultado numérico da simulação. O primeiro ponto a observar é que a queda de $\pi$ reduz diretamente a força do primeiro estágio. Esse comportamento é ilustrado pela Figura~\ref{fig:q14_f_primeiro_estagio}, que mostra a relação entre o coeficiente do primeiro estágio e o $F$ médio da regressão de primeiro estágio. Quando $\pi$ é elevado, o instrumento explica bem a variável endógena e o $F$ médio fica muito acima do limiar convencional de 10. Quando $\pi$ diminui, o $F$ médio cai rapidamente, indicando que o instrumento perde relevância estatística.

\begin{figure}[H]
    \centering
    \includegraphics[width=0.88\textwidth]{Lista_1/Figures/Stata/stata_fig08_simulation_ffs.pdf}
    \caption{Força do primeiro estágio por nível de $\pi$ na simulação de instrumentos fracos. Cada ponto mostra a estatística $F$ média do primeiro estágio em 250 repetições. A linha vermelha indica o limiar convencional $F=10$.}
    \label{fig:q14_f_primeiro_estagio}
\end{figure}

Os resultados mostram que, quando o instrumento é forte, isto é, quando $\pi=1$, o estimador 2SLS se concentra em torno do valor verdadeiro. A média estimada foi aproximadamente $-0,998$, muito próxima de $\beta_p=-1$, com viés praticamente nulo e desvio-padrão de apenas $0,064$. Além disso, como mostra a Figura~\ref{fig:q14_f_primeiro_estagio}, o $F$ médio do primeiro estágio foi bastante elevado, aproximadamente $208,87$, indicando forte capacidade explicativa do instrumento sobre a variável endógena.

À medida que $\pi$ diminui, o primeiro estágio se enfraquece. Para $\pi=0,25$, o $F$ médio cai para aproximadamente $13,40$, ainda próximo da regra de bolso usual, mas já acompanhado de aumento do viés e da dispersão do estimador. A média do 2SLS passa para $-1,100$ e o desvio-padrão sobe para $0,392$. Esse resultado já indica que a perda de força do instrumento torna a estimação menos precisa.

O problema se torna mais evidente quando $\pi$ cai para $0,10$, $0,05$ e $0,02$. Nesses casos, o $F$ médio do primeiro estágio fica abaixo de 10, chegando a aproximadamente $3,35$, $1,66$ e $1,11$, respectivamente. O estimador 2SLS passa a apresentar comportamento extremamente instável. Para $\pi=0,10$, a média estimada é $-0,479$, afastando-se bastante do valor verdadeiro. Para $\pi=0,05$, a média cai para $-4,596$, com desvio-padrão de $58,596$. Para $\pi=0,02$, a média muda inclusive de sinal, chegando a $1,130$, com desvio-padrão de $48,751$.

A Figura~\ref{fig:q14_simulacao_instrumentos_fracos} sintetiza visualmente esse resultado. O painel A compara a estimativa média por 2SLS, a estimativa média por MQO e o valor verdadeiro do parâmetro estrutural. Ele mostra que, quando o instrumento é forte, a estimativa por 2SLS fica próxima de $\beta_p=-1$. Quando o instrumento enfraquece, a estimativa se torna instável e pode se afastar substancialmente do valor verdadeiro. Esse comportamento ilustra o problema de viés em amostras finitas associado a instrumentos fracos.

O painel B da Figura~\ref{fig:q14_simulacao_instrumentos_fracos} mostra diretamente o viés absoluto do 2SLS. A mensagem é clara: quando o instrumento é forte, o viés é pequeno; quando o instrumento perde relevância, o viés aumenta de forma expressiva. O painel C complementa essa interpretação ao mostrar a dispersão do estimador. A elevação do desvio-padrão quando $\pi$ é baixo indica que o 2SLS não apenas se torna mais enviesado, mas também muito menos preciso.

Por fim, o painel D da Figura~\ref{fig:q14_simulacao_instrumentos_fracos} mostra a cobertura dos intervalos convencionais de 95\%. A cobertura não se deteriora de forma perfeitamente monotônica na simulação: cai para 92,4\% quando $\pi=0,05$, mas volta a 97,2\% quando $\pi=0,02$. Esse aparente retorno da cobertura não significa que a inferência convencional voltou a ser confiável. Pelo contrário, quando o instrumento é extremamente fraco, os intervalos podem se tornar muito largos e instáveis, cobrindo o valor verdadeiro por excesso de imprecisão, não por boa qualidade inferencial. Assim, a cobertura isolada pode esconder o fato central: a estimativa perdeu precisão e conteúdo econômico.

\begin{figure}[H]
    \centering
    \includegraphics[width=1.1\textwidth]{Lista_1/Figures/Stata/stata_fig09_q14_simulacao_instrumentos_fracos.pdf}
    \caption{Simulação de instrumentos fracos e desempenho do estimador 2SLS. A força do instrumento é medida pelo coeficiente do primeiro estágio, $\pi$. Quando $\pi$ diminui, o instrumento perde capacidade de explicar a variável endógena. A figura mostra que a redução da força do instrumento aumenta o viés e a dispersão do estimador 2SLS, além de tornar a cobertura dos intervalos convencionais de 95\% menos informativa.}
    \label{fig:q14_simulacao_instrumentos_fracos}
\end{figure}

Portanto, a simulação confirma a intuição discutida nas questões anteriores: quando o instrumento perde relevância, o 2SLS deixa de se comportar como um estimador bem concentrado em torno do parâmetro verdadeiro. O viés aumenta, a dispersão cresce fortemente e os resultados passam a depender muito mais da variação amostral. Isso ajuda a interpretar os problemas encontrados nas especificações empíricas anteriores, especialmente nos modelos em que o $F$ efetivo ou o $F$ usual do primeiro estágio indicaram instrumentos fracos.

Esse resultado dialoga diretamente com a discussão sobre instrumentos fracos feita anteriormente. Na presença de instrumentos pouco relevantes, o estimador de variáveis instrumentais pode se tornar viesado, altamente disperso e pouco confiável para inferência. Por isso, em aplicações de Economia Industrial, sobretudo quando a elasticidade-preço estimada será usada para discutir poder de mercado, margem ou perda crítica, é perigoso interpretar mecanicamente o coeficiente 2SLS sem antes avaliar a força do primeiro estágio.

\end{enumerate}

\end{document}